\documentclass[11pt,reqno]{amsart}
\usepackage[T1]{fontenc}
\usepackage{lmodern}
\usepackage[margin=1in]{geometry}
\usepackage{amsmath,amssymb,mathtools,booktabs,microtype}
\usepackage[hidelinks]{hyperref}
\hypersetup{
 pdftitle={Clipped Gram energies and zero counts of the Riemann zeta function},
 pdfauthor={R. Arun Chandru},
 pdfsubject={Unconditional zero-counting bounds, phase rigidity, and variational method ceilings},
 pdfkeywords={Riemann zeta function, simple zeros, critical line, distinct zeros, Gram certificates, phase rigidity}
}
\numberwithin{equation}{section}
\newcommand{\R}{\mathbb R}
\newcommand{\C}{\mathbb C}
\newcommand{\Z}{\mathbb Z}
\newcommand{\tr}{\operatorname{tr}}
\newcommand{\dist}{\operatorname{dist}}
\newcommand{\sinc}{\operatorname{sinc}}
\newcommand{\supp}{\operatorname{supp}}
\newcommand{\osc}{\operatorname{osc}}
\newcommand{\dd}{\,\mathrm d}
\newcommand{\Ns}{N_{\rm s}}
\newcommand{\Nsc}{N_{\rm s0}}
\newcommand{\Nd}{N_{\rm d}}
\newcommand{\Nc}{N_0}
\newcommand{\one}{\mathbf1}
\newcommand{\qstar}{q_\star}
\title[Clipped Gram energies and zero counts]{Clipped Gram energies and zero counts \\of the Riemann zeta function}
\author{R. Arun Chandru}
\email{arun.chandru@gmail.com}
\date{9 October 2026}
\subjclass[2020]{Primary 11M26; Secondary 15A42, 42A16, 65G20}
\keywords{Riemann zeta function, simple zeros, Gram matrices, phase rigidity, variational bounds, rational certificates}
\begin{document}
\raggedbottom
\begin{abstract}
We study the interaction energy that remains available when the Gram
matrix associated with the zeros of the Riemann zeta function is
spectrally clipped. An explicit band-limited majorant and a close-pair
estimate transfer an existing eight-point inequality to the full clipped
energy, uniformly over configurations of arbitrary size. Together with
the unconditional pair-correlation formula this gives
\[
 \liminf_{T\to\infty}\frac{N_{\rm s0}(T)}{N(T)}
 \ge\frac{66812491}{99194876}=0.673547805029767868\ldots,
\]
where \(N_{\rm s0}\) counts simple critical-line zeros and \(N\) counts
all nontrivial zeros with multiplicity, giving \(67.35478050\%\). The corresponding average of
the simple-zero and critical-line proportions is at least
\(83.67739025\%\). Evaluating separate labelled and multiplicity
certificates gives \(83.90046180\%\) for distinct locations,
\(67.80092360\%\) for simple zeros anywhere,
\(88.93165002\%\) for simple-or-critical-line mass, and
\(83.79775462\%\) for mass at locations of multiplicity at most two;
such locations constitute at least \(93.94524765\%\) of the distinct
locations. All terminating lower percentages are rounded down.
Beyond these numerical deductions, we prove a phase comparison in which
local slack pays continuously for long physical spans, derive a common
residual for joint profile-and-threshold deductions, and solve a
two-lattice variational problem explicitly. A twelve-fibre construction
extends a periodic-mixture obstruction to the full affine \(L^2\) space
and gives a \(67.404\%\) ceiling for the specified joint-affine
simple-line relaxation. A simultaneous union ceiling is
\(89.497833220444\%\). These ceilings describe limitations of the
relaxation, not upper bounds on the actual zero proportions.
\end{abstract}
\maketitle
\tableofcontents

\section{Introduction}

A multiple zero is expensive in a quadratic energy: multiplicity \(j\)
contributes \(j\) to the zero count but \(j^2\) to the diagonal
energy. Nearby distinct zeros also interact. A lower bound for this
interaction leaves less of the available energy for multiplicities and
for zeros away from the critical line. The central issue is to retain
that interaction through the spectral inequality which makes the
argument unconditional.

For a nonnegative even density \(f\) of integral one on
\(I=[-1/2,1/2]\), let \(K_f\) denote its Fourier transform. On the
critical line, normalized zero locations \(y_i\) determine unit vectors
whose Gram matrix has entries \(K_f(y_i-y_j)\). Its off-diagonal
square energy measures the interaction. Off the critical line,
reflection symmetry gives Hermitian terms with at most one positive
direction per reflected pair. It does not give positive rank-one
terms. The appropriate comparison therefore uses a clipped square
\[
 \phi_c(x)=x^2-(x-c)_+^2
\]
rather than the full square on an arbitrary retained Gram matrix.
Recovering full energy where it is legitimate, and keeping track of
what is lost elsewhere, will be the organizing principle of the paper.

The analytic foundation is the unconditional Montgomery theorem of
Baluyot, Goldston, Suriajaya and Turnage-Butterbaugh
\cite{BGSTB}. It evaluates a pair sum over the complex zeros without
first placing them on the critical line. Their subsequent work
\cite{BGSTBI} studies the horizontal information available from pair
correlation under a narrow-strip hypothesis. The unconditional
finite-rank and Hilbert-space arguments of Alp\"oge--Furman and
Lamzouri \cite{AF,Lamzouri} handle the off-line pairs by their inertia
instead. With the Montgomery--Taylor profile they yield more than
\(67.25\%\) simple critical-line zeros. These arguments descend from
Montgomery's pair-correlation method and its test-function
optimization \cite{Montgomery,MontgomeryTaylor}; conditional
semidefinite refinements appear in \cite{CGdL}. The present argument
uses the unconditional complex-zero formula throughout.

Local Gram inequalities have sharpened the basic diagonal estimate.
The relevant developments include overlapping finite windows,
pressure optimization, bounded correction potentials and compatible
window systems
\cite{Ainta,AMTOPA,Tawanerguo,ChandruP6,GebendorferWindows,GebendorferExpanded}.
Close-pair separation allows the full energy of a separated remainder
to survive clipping \cite{Freund,DSouza,Knausgard2026}. Labelled
inequalities distinguish simple and double locations
\cite{KnausgardDistinct,Santibanez,Knausgard2026}, while residual and
multiplicity accounting support further consequences
\cite{Lamzouri,ChandruP6,ShiTrace,Rakitin,Waddle}.
Related numerical refinements and their public certificate data are
reported in \cite{OjhaTrmdy,Schafer,Singh}.
These ingredients have different roles. A stronger finite inequality
is useful only when its Fourier profile, its trace energy and the
clipped operator to which it is applied remain the same.

Our principal numerical result is
\begin{equation}
 \liminf_{T\to\infty}\frac{\Nsc(T)}{N(T)}
 \ge q_*:=\frac{66812491}{99194876}
 =0.673547805029767868\ldots .
 \label{eq:intro-main}
\end{equation}
Here \(N(T)\) counts nontrivial zeros with \(0<\Im\rho\le T\),
with analytic multiplicity, and \(\Nsc(T)\) counts those that are
simple and lie on \(\Re\rho=1/2\). The numerical profile and its
eight-point floor are specified public inputs \cite{AMTOPA,Lavery2026}.
For that profile we give an elementary explicit majorant with Fourier
support \([-17/20,17/20]\) and integral strictly below two, together
with a close-pair kernel lower bound. Every separated Gram matrix is
then below the clipping level, regardless of its size. A maximal
matching of close pairs handles the complement. Summing the local
inequality on the separated remainder uses the full lag capacities
and incurs only a fixed endpoint loss. The smoothing needed by the
analytic formula costs \(O(\delta n)\), not \(O(\delta n^2)\).
These facts turn the local floor into the global estimate
\[
 \tr\phi_2(U)-n\ge
 \frac{805124}{10^8}n-\frac{404350}{10^8}S-O(1),
\]
where \(S\) is the span and \(U\) is the normalized Gram matrix.
This realization of an established separation mechanism, for the
specified profile with explicit uniform constants, yields
\eqref{eq:intro-main}.

Keeping the positive-index information also gives
\[
 \liminf_{T\to\infty}\frac{\Ns(T)+\Nc(T)}{2N(T)}
 \ge\frac{1+q_*}{2}=0.836773902514883934\ldots .
\]
The two counts in this average have different meanings:
\(\Ns\) counts simple zeros anywhere, whereas \(\Nc\) counts all
critical-line multiplicity. They are not each bounded below by the
displayed average. Table~\ref{tab:plus-counts} collects seven
consequences, with every numerator and denominator identified.
The remaining five values come from evaluating other stated
certificates at sharper exact energy bounds, not from adding the
gains of incompatible profiles. The small numerical refinements
are kept distinct from the structural results.

The first structural question concerns information left after a local
inequality is summed. Its nonnegative slack \(H\) and the energy \(T\)
at unused index lags are often both discarded. A sinc-type lower bound
for the Fourier kernel converts small long-lag energy into phase
agreement. The difficulty is that a fixed index lag can have an
arbitrarily large physical span. We show that the sum of the local
slacks over a frame cover pays for that span continuously. The result
is a uniform estimate for the squared distance of adjacent gaps to
the integers. For one explicit profile,
\begin{equation}
 \sum_i\dist(y_{i+1}-y_i,\Z)^2\le350(H+T)+6.
 \label{eq:intro-rigidity}
\end{equation}
More generally, under fixed local slack and endpoint hypotheses, a
degree-\(D\) cosine profile permits an \(O(D)\) energy price. The
argument keeps a geometric residual which an affine count alone
does not record. Its profile is separate from that of
\eqref{eq:intro-main}.

The complementary structural question is how much any deduction from
a specified collection of affine floors can establish. Periodic
configurations supply tests of every universal floor. We put their
profile dependence, the clipping thresholds and the unused test
slack into one exact residual. The residual permits arbitrary fixed
finite combinations of profiles, with separate thresholds, while
testing all of them on common configurations. It leads to two
variational conclusions. A mixture of the unit lattice and the
spacing-two lattice has an explicit positive minimizer, a coercive
error term and an exact two-count frontier within that
two-test relaxation. A period-three/period-four mixture, originating
in the finite-dimensional obstruction of
\cite{GebendorferCap}, admits a twelve-fibre formulation on all of
\(L^2(I)\). A rational approximate minimizer and a rigorous residual
bound yield the \(67.404\%\) ceiling stated in the abstract.
The two-lattice frontier also gives the simultaneous
\(89.497833220444\%\) union ceiling.

A ceiling here is a feasibility statement about retained inequalities.
It says that those inequalities allow a particular abstract count
population; it does not construct zeros with those counts. This
distinction is productive: it identifies the extra information a
stronger deduction must retain. The phase comparison retains such
information, as do labelled-block constraints not included in the
specified relaxation. In a smaller scalar multiplicity relaxation we
can be more explicit: a population of simple and triple real
locations satisfies every scalar clipping row and simultaneously
attains the two cutoff-two bounds. Changing the clipping level alone
cannot improve these bounds.

We first establish the signed-operator interface and prove the
majorant and full-energy transfer for the main profile. We then give
the count deductions, the phase comparison, and the joint
variational analysis. The appendices specify the finite rational
calculations and the data used in the twelve-fibre estimate.
Computer-assisted local inequalities are cited with their precise
profiles and domains; all subsequent operator, counting and
variational arguments are given here. The accompanying formal source
covers the main numerical simple-line chain, while the additional
count and variational results have the proofs and arithmetic
certificates specified in the paper.
\section{One operator and the counting interface}
\label{sec:interface}

\subsection{Profiles, complex zeros, and the trace energy}

Put $I=[-1/2,1/2]$. For an even probability density $f\in L^2(I)$,
extended by zero, define
\begin{equation}
 K_f(t)=\int_I f(u)e^{2\pi itu}\dd u,\qquad
 R(f)=\int_I f^2+\iint_{I^2}|s-t|f(s)f(t)\dd s\dd t.
 \label{eq:normalization}
\end{equation}
The same integral defines $K_f(t)$ for complex $t$ and is entire,
because $f$ is integrable with compact support.
We use $\sinc(t)=\sin(\pi t)/(\pi t)$, with $\sinc(0)=1$.
For the analytic construction in this section assume additionally that
$\sqrt f$ is smooth in the interior of $I$. The numerical profiles are
strictly positive finite cosine sums and satisfy this assumption. The
variational problems later in the paper are defined on all of $L^2(I)$.
Choose smooth even cutoffs $0\le\chi_\varepsilon\le1$, supported inside
$I$, tending to one, and set
\[
 m_\varepsilon=\int f\chi_\varepsilon^2,\qquad
 f_\varepsilon=f\chi_\varepsilon^2/m_\varepsilon
 =\eta_\varepsilon^2.
\]
Then $f_\varepsilon\to f$ in $L^1$ and $L^2$, and
\begin{equation}
 K_\varepsilon(0)=1,\quad |K_\varepsilon|\le1,\quad
 \sup_{t\in\R}|K_\varepsilon(t)-K_f(t)|
 \le\|f_\varepsilon-f\|_1\longrightarrow0.
 \label{eq:smoothing}
\end{equation}

Let $N(T)$ count all nontrivial zeros $\rho$ with
$0<\Im\rho\le T$, with analytic multiplicity. Write $L=\log T$ and
$z_\rho=i(\rho-1/2)L/(2\pi)$. Reflection
$\rho\mapsto1-\bar\rho$ becomes $z\mapsto\bar z$.
On the finite-dimensional span of
$F_z(u)=\eta_\varepsilon(u)e^{-2\pi iuz}$ consider
\begin{equation}
 \mathcal A=\sum_z\nu_zF_z\langle F_{\bar z},\cdot\rangle,
 \qquad
 \tr\mathcal A=N(T),\quad
 \tr\mathcal A^2=\sum_{z,w}\nu_z\nu_w K_\varepsilon(z-w)^2.
 \label{eq:operator}
\end{equation}
The inner product is conjugate-linear in its first argument. Conjugate
pairs make $\mathcal A$ Hermitian. The square in the last expression is
not a modulus square; the whole sum is nevertheless the nonnegative
trace of the square of a Hermitian operator.

The unconditional analytic input, in the form of
\cite[Lemma 3.1 and Section 3]{Lamzouri} and
\cite[Proposition 3.1]{Knausgard2026}, is
\begin{equation}
 \tr\mathcal A^2=(R(f_\varepsilon)+o(1))N(T),\qquad
 R(f_\varepsilon)\longrightarrow R(f).
 \label{eq:trace-limit}
\end{equation}
The first limit keeps $\varepsilon$ fixed. To specify the normalization,
put $g=f_\varepsilon*f_\varepsilon$. Then
\[
 \widehat{g-g''/(4L^2)}(x)
 =(1+\pi^2x^2/L^2)K_\varepsilon(x)^2.
\]
At $x=z_\rho-z_{\rho'}$ the factor cancels the weight
$4/(4-(\rho-\rho')^2)$ in the unconditional pair formula
\cite{BGSTB}. Both $g$ and $g''$ are fixed admissible smooth tests;
the derivative contribution is $O_\varepsilon(N/L^2)$, while the first
contributes $g(0)+\int|u|g(u)\dd u=R(f_\varepsilon)$.
The scaled real coordinates have total span at most
$T\log T/(2\pi)=N(T)+o(N(T))$.

\subsection{The finite inequality and simple critical-line zeros}

For $c\ge0$ let
\begin{equation}
 \phi_c(x)=x^2-(x-c)_+^2.
 \label{eq:clipping}
\end{equation}
It is convex, quadratic below $c$, and linear above $c$. If $A=P+Q$
is Hermitian, $P\succeq0$, and $Q$ has at most $p$ positive eigenvalues,
then for $t\ge c\ge0$,
\begin{equation}
 \tr A^2\ge\tr\phi_c(P)+2t(\tr A-\tr P)-t^2p.
 \label{eq:two-level}
\end{equation}
Indeed write $Q=Q_+-Z$, where $Z\succeq0$ and $ZQ_+=0$.
The cross term gives
$\tr A^2\ge\tr(P-Z)^2+\tr Q_+^2$.
For $X=P-cI$, the inequalities
$\tr(XZ)\le\tr(X_+Z)$ and $\tr(X_+-Z)^2\ge0$ imply
$\tr(P-Z)^2+2c\tr Z\ge\tr\phi_c(P)$.
Also $\tr Q_+^2\ge2t\tr Q_+-t^2p$.
Adding leaves $2(t-c)\tr Z\ge0$, proving
\eqref{eq:two-level}. This is the finite clipping argument of
\cite{KnausgardDistinct,Knausgard2026}, recalled to fix every coefficient.

For any principal-block partition of a Hermitian matrix $G$,
\begin{equation}
 \tr\phi_c(G)\ge\sum_B\tr\phi_c(G_B).
 \label{eq:pinching}
\end{equation}
Choose an orthonormal eigenbasis of each block and apply scalar Jensen
to the spectral measure of $G$ at each resulting basis vector. Summing
proves the assertion; operator convexity is not needed.

Let $n=\Nsc(T)$ be the number of simple critical-line zeros, and take
$P$ to be their part of \eqref{eq:operator}. Its nonzero spectrum is
that of the Gram matrix $U=(K_\varepsilon(y_i-y_j))$, and $\tr P=n$.
Each remaining real location contributes at most one positive direction
and has multiplicity at least two. Each reflected off-line pair of
multiplicity $j$ contributes $j(vw^*+wv^*)$, of trace $2j$ and with
at most one positive direction, since
\[
 vw^*+wv^*=\tfrac12\bigl((v+w)(v+w)^*-(v-w)(v-w)^*\bigr).
\]
Thus $2p\le N-n$. With $c=t=2$,
\begin{equation}
 \tr\mathcal A^2\ge2(N-n)+\tr\phi_2(U).
 \label{eq:simple-master}
\end{equation}

Suppose that $a<1$, $b\ge0$, and a constant $B$ satisfy, for every
finite real configuration of $n$ points and all small $\varepsilon$,
\begin{equation}
 \tr\phi_2(U)-n\ge an-bS-\omega(\varepsilon)n-B,
 \qquad\omega(\varepsilon)\longrightarrow0,
 \label{eq:affine-certificate}
\end{equation}
where $S$ is its span (zero for at most one point). Combining
\eqref{eq:trace-limit}--\eqref{eq:affine-certificate}, and taking
$T\to\infty$ before $\varepsilon\to0$, gives
\begin{equation}
 \liminf_{T\to\infty}\frac{\Nsc(T)}{N(T)}
 \ge\frac{2-R(f)-b}{1-a}.
 \label{eq:simple-quotient}
\end{equation}
This is the standard certificate interface
\cite{AF,Lamzouri,ChandruP6,Knausgard2026}. The numerical application below keeps its profile, finite local input,
and global Gram operator compatible throughout.

\section{A profile for which separated energy is retained}
\label{sec:am-profile}

The numerical construction uses the interval $I=[-1/2,1/2]$ and the
positive profile
\begin{equation}
 v(s)=\cos(\sqrt2s)+\sum_{j=1}^{12}d_j\cos(2\pi js),\qquad
 f(s)=\frac{v(s)}{K_0},\qquad
 K_0=\sqrt2\sin\theta,\quad \theta=\frac{\sqrt2}{2}.
 \label{eq:am-profile}
\end{equation}
The density $f$ is extended by zero outside $I$. The coefficients are
specified exactly by
\[
\begin{array}{c|r@{\qquad}c|r}
j&10^9d_j&j&10^9d_j\\\hline
1&12310798&7&-6846472\\
2&-15041681&8&3781297\\
3&3867664&9&-5670353\\
4&6489926&10&3355089\\
5&-992327&11&-450523\\
6&5580097&12&-218483
\end{array}
\]
This is the cosine family introduced in \cite{AMTOPA}, with the
coefficients and eight-point inequality of \cite{Lavery2026}.
We apply that inequality to the full clipped Gram energy. Close-pair matching, separated-set norm
bounds, and convex pinching are established parts of this method
\cite{Freund,Knausgard2026}; the estimates here realize their requirements
for the specific profile \eqref{eq:am-profile}. In fact, its coefficient
absolute sum also lies within the uniform cosine-family range of
\cite[Section 3]{Freund}. Thus the proportion obtained below is a new
evaluation of the established transfer, rather than a new general
separation theorem. The explicit estimates given here provide a
self-contained realization at separation $17/20$.

The normalization is immediate because every correction cosine has
integral zero. Moreover,
\[
 \sum_{j=1}^{12}|d_j|\le\frac{13}{200},\qquad
 \cos(\sqrt2s)\ge1-s^2,
\]
so $v(s)\ge137/200$ on $I$. In particular $f$ is a probability
density, and $|K_f(t)|\le1$ for real $t$.

Its energy has a particularly simple exact expression. Set
\begin{equation}
 P=\sum_{j=1}^{12}d_j^2
       \left(\frac12-\frac1{4\pi^2j^2}\right),\qquad
 H=2-R(f)
  =\frac32-\frac{\cos\theta}{\sqrt2\sin\theta}
       -\frac{P}{2\sin^2\theta}.
 \label{eq:am-energy}
\end{equation}
To see the cancellation behind this formula, put
$v_0(s)=\cos(\sqrt2s)$. The function
\[
 v_0(s)+\int_I|s-t|v_0(t)\,dt
\]
has zero second derivative and zero derivative at $s=1/2$; it is
therefore the constant $\cos\theta+K_0/2$. All cross terms with a
mean-zero perturbation vanish in the energy form. For $j\ge1$, direct
differentiation gives
\[
 \int_I|s-t|\cos(2\pi jt)\,dt
   =\frac{(-1)^j-\cos(2\pi js)}{2\pi^2j^2}.
\]
Orthogonality of the cosines then gives the correction $P$, while
the unperturbed energy is $K_0\cos\theta+K_0^2/2$. Division by
$K_0^2=2\sin^2\theta$ proves \eqref{eq:am-energy}.

Here and below the numerical enclosures can be checked with rational
arithmetic. One convenient enclosure is
\begin{equation}
 H\ge H_0:=\frac{67216841}{10^8}.
 \label{eq:am-H-floor}
\end{equation}
For completeness, introduce the rational numbers
\[
\begin{aligned}
 r_-&=\frac{7071067811865475244}{10^{19}},&
 r_+&=\frac{7071067811865475245}{10^{19}},\\
 s_-&=\frac{64963693904}{10^{11}},&
 c_+&=\frac{76024459708}{10^{11}},\\
 p_+&=\frac{314159265358979323847}{10^{20}}.
\end{aligned}
\]
Squaring brackets $\theta$ between $r_-$ and $r_+$. Taylor bounds give
$\sin\theta\ge s_-$ and $\cos\theta\le c_+$, and $\pi<p_+$.
For example, the sine and cosine polynomials of degree at most eleven
have error at most $13|x|^{12}/5748019200$ for $|x|\le1$; applying
them at the rational endpoints proves the two trigonometric bounds.
With
\[
 P_+=\sum_{j=1}^{12}d_j^2\left(\frac12-\frac1{4p_+^2j^2}\right),
\]
the remaining check is the exact rational inequality
\begin{equation}
 \frac32-\frac{c_+}{2r_-s_-}-\frac{P_+}{2s_-^2}
 \ge\frac{67216841}{10^8}.
 \label{eq:am-H-rational-check}
\end{equation}
The positive denominators and $0\le P\le P_+$ give the direction
required in \eqref{eq:am-H-floor}.

\subsection{A uniform norm bound from separation}

Let $d=17/20$ and use the normalization
$\operatorname{sinc}(t)=\sin(\pi t)/(\pi t)$. Consider
\begin{equation}
 g(s)=\frac7{10}\left\{
  \operatorname{sinc}\bigl(d(s-1/2)\bigr)+
  \operatorname{sinc}\bigl(d(s+1/2)\bigr)\right\}^2.
 \label{eq:am-majorant}
\end{equation}
This nonnegative integrable function satisfies
\begin{equation}
 f(s)+\frac{5789}{960000}\le g(s)\quad(s\in I),\qquad
 \Lambda:=\int_{\mathbb R}g(s)\,ds\le\frac{1253}{650}<2,
 \qquad \widehat g(t)=0\quad(|t|\ge d).
 \label{eq:am-majorant-properties}
\end{equation}
The strict margin and the Fourier support, rather than pointwise
convergence of matrix entries, are what make the resulting norm bound
independent of the number of points.

The pointwise inequality admits a short polynomial verification.
Since $K_0^{-1}\le109/100$, the elementary upper bound for cosine gives,
with $x=s^2\in[0,1/4]$,
\[
 f(s)\le\frac{109}{100}
       \left(\frac{213}{200}-x+\frac{x^2}{6}\right).
\]
To bound the sum of sincs from below, write $D=\pi d$ and use
\[
 D^2\le\frac{357}{50},\qquad
 D^4\ge\frac{31684}{625},\qquad D^6\le364.
\]
These follow from $157/50\le\pi\le22/7$. The degree-six lower
bound for $\sin z/z$, applied at $D(s-1/2)$ and $D(s+1/2)$,
gives the polynomial
\[
\begin{split}
 &2-\frac{119}{200}+\frac{7921}{150000}-\frac{13}{5760}
 +\left(-\frac{119}{50}+\frac{7921}{6250}-\frac{13}{96}\right)x\\
 &\hspace{35mm}
 +\left(\frac{31684}{37500}-\frac{13}{24}\right)x^2
 -\frac{13}{90}x^3.
\end{split}
\]
Its constant exceeds $29/20$, its linear coefficient exceeds $-5/4$,
and its last two terms have nonnegative sum on $[0,1/4]$.
Consequently the sinc sum is at least $29/20-5x/4>0$.
Subtracting the displayed upper bound for $f$ from the corresponding
lower bound for $g$ leaves
\[
 Q(x)=\frac{3109}{10000}-\frac{579}{400}x+
                         \frac{2189}{2400}x^2.
\]
Here
\[
 Q(x)-Q(1/4)
  =(1/4-x)\left\{\frac{579}{400}
                   -\frac{2189}{2400}(x+1/4)\right\}\ge0,
 \qquad Q(1/4)=\frac{5789}{960000},
\]
which proves the required margin. The degree-seven lower Taylor
bound for sine also gives
\[
 K_0\ge1-\frac1{12}+\frac1{480}-\frac1{40320}
       >\frac{100}{109},
\]
and hence the bound $K_0^{-1}\le109/100$ used above.

The Fourier transform of each sinc in \eqref{eq:am-majorant} is a
modulated interval indicator. The square corresponds to convolution,
so its transform is supported on $[-d,d]$, is continuous, and vanishes
at both endpoints. Parseval gives
\[
 \Lambda=\frac{28}{17}\{1+\operatorname{sinc}(17/20)\}.
\]
For a rational bound on the last quantity, put $y=3\pi/20$ and use
$\sin(17\pi/20)=\sin y$. The estimates
$y^2\ge22206/10^5$ and $y^4\le4932/10^5$ imply
\[
 \operatorname{sinc}(17/20)
 \le\frac3{17}\left(1-\frac{y^2}{6}+\frac{y^4}{120}\right)
 \le\frac3{17}\frac{963401}{10^6}.
\]
Substitution proves the bound $\Lambda\le1253/650$.

For any $d$-separated real points $x_1,\ldots,x_m$ and complex
coefficients $z_1,\ldots,z_m$, \eqref{eq:am-majorant-properties}
therefore gives
\begin{equation}
\begin{split}
 \sum_{i,j}\overline{z_i}z_jK_f(x_i-x_j)
 &=\int_I f(s)\left|\sum_jz_je^{-2\pi ix_js}\right|^2\,ds\\
 &\le\int_{\mathbb R}g(s)
                  \left|\sum_jz_je^{-2\pi ix_js}\right|^2\,ds
 =\Lambda\sum_j|z_j|^2.
\end{split}
\label{eq:am-separated-norm}
\end{equation}
The identity uses the vanishing of the off-diagonal Fourier terms,
including pairs at distance exactly $d$. Thus the Gram matrix on
every such set has spectrum below two.

The same conclusion applies after the smoothing used in the trace
formula. Indeed $m_\varepsilon f_\varepsilon=f\chi_\varepsilon^2\le f$
and $m_\varepsilon\to1$, so the smoothed Gram matrix has norm at
most $\Lambda/m_\varepsilon<2$ for all sufficiently small
$\varepsilon$. This choice of $\varepsilon$ works simultaneously for
all separated sets and all their cardinalities.

\subsection{What a close pair contributes}

The complementary estimate is
\begin{equation}
 K_f(t)\ge\frac{2583}{20000}>\frac18
 \quad (|t|\le17/20).
 \label{eq:am-close-pair}
\end{equation}
Here is a direct bound that also keeps removable singularities out of
the argument. Put
\[
 K_{\rm MT}(t)=\int_I\cos(\sqrt2s)\cos(2\pi ts)\,ds.
\]
For $0\le t\le d$, this function decreases with $t$, because the
cosine factor is nonnegative and the other cosine decreases pointwise
when expressed using $|s|$. At the endpoint, elementary integration
gives, with $z=17\pi/20$, $y=3\pi/20$, and
$\gamma=\cos\theta/K_0$,
\[
 \frac{K_{\rm MT}(d)}{K_0}
  =\frac{\cos y+2\gamma z\sin y}{2z^2-1}\ge\frac15.
\]
For the last comparison it suffices to use
\[
 \gamma\ge\frac45,\quad \frac83\le z\le\frac{187}{70},
 \quad\sin y\ge\frac9{20},\quad\cos y\ge\frac{22}{25}.
\]
The numerator is then at least $14/5$ and the positive denominator at
most $32519/2450$. The sine and cosine bounds follow from their first
Taylor inequalities and $471/1000\le y\le33/70$.
For $\gamma$, use $\cos\theta\ge3/4$ and
$K_0\le1-1/12+1/480=147/160$, which give
$\gamma\ge40/49>4/5$.

The transform of each correction cosine is
\[
 \frac12\{\operatorname{sinc}(t-j)+
                  \operatorname{sinc}(t+j)\},
\]
of absolute value at most one. The total correction to
$K_{\rm MT}/K_0$ is therefore at most
\[
 (109/100)(13/200)=1417/20000.
\]
Combining this with the preceding
endpoint estimate proves \eqref{eq:am-close-pair}; evenness covers
negative $t$.

\section{The eight-point floor and its global consequence}
\label{sec:am-transfer}

Take nonnegative gaps $g_0,\ldots,g_6$, and put
$p_0=0$, $p_j=\sum_{r<j}g_r$ for $1\le j\le7$.
The prior eight-point certificate \cite{Lavery2026} is the inequality
\begin{equation}
 \sum_{r=0}^6b_rg_r+
 \sum_{0\le i<j\le7}a_{ij}K_f(p_j-p_i)^2
 \ge a_8:=\frac{805124}{10^8}
 \qquad(g_r\ge0).
 \label{eq:am-local-floor}
\end{equation}
All entries in the following table are the numerators of $10^8a_{ij}$;
the lower triangle is unused.
\begin{center}
\small
\setlength{\tabcolsep}{3.5pt}
\begin{tabular}{c|rrrrrrr}
$i\backslash j$&1&2&3&4&5&6&7\\\hline
0&13630830&0&19565904&20317441&45030671&100000000&200000000\\
1&&29997996&30621878&47766489&79682558&109938656&100000000\\
2&&&37356140&69378121&65335210&79682558&45030671\\
3&&&&38030065&69378121&47766489&20317441\\
4&&&&&37356140&30621878&19565904\\
5&&&&&&29997996&0\\
6&&&&&&&13630830
\end{tabular}
\end{center}
The pressure coefficients are
\begin{equation}
 10^8(b_0,\ldots,b_6)
 =(28898,57272,75526,80958,75526,57272,28898),\qquad
 B_8:=\sum_rb_r=\frac{404350}{10^8}.
 \label{eq:am-pressures}
\end{equation}
The coefficients are nonnegative. Their seven lag sums are
\begin{equation}
\begin{split}
10^8\left(\sum_{i=0}^{7-\ell}a_{i,i+\ell}\right)_{\ell=1}^7
&=(199999997,199999998,199999996,\\
&\quad199999998,199999998,200000000,200000000),
\end{split}
\label{eq:am-capacities}
\end{equation}
so each lag has capacity at most two.

The local floor \eqref{eq:am-local-floor} is the computer-assisted
part of this input. Its domain and every coefficient are specified
here, while the accompanying formal source contains its complete
finite proof data. Since every $b_r$ is positive, a gap with
$b_rg_r\ge a_8$ is handled by pressure alone. On the remaining
bounded box the verification uses rational cells with retained
pair-distance constraints, proved analytic minorants for $K_f^2$,
and nonnegative multiplier combinations. Infeasible cells are excluded
by rational linear combinations of their constraints; the other leaves
verify the stated floor. The coefficient table is invariant under
reversing the seven gaps, so a certified half-domain suffices.
The analytic enclosures, covering argument, and exact leaf data are
included in that source. Thus the numerical floor is a specified
reusable mathematical input; the finite argument that follows only
requires \eqref{eq:am-local-floor} and the visible capacities
\eqref{eq:am-capacities}.

Consider first ordered points $x_0<\cdots<x_{m-1}$ with $m\ge8$.
Sum \eqref{eq:am-local-floor} over all consecutive eight-point
windows. If $S=x_{m-1}-x_0$, the pressure telescopes exactly to
\begin{equation}
 \sum_{r=0}^6 b_r(x_{m-7+r}-x_r)\le B_8S.
 \label{eq:am-pressure-telescope}
\end{equation}
Each pair of index lag $\ell\le7$ is charged at most the full
lag sum in \eqref{eq:am-capacities}. Hence
\begin{equation}
 E_7(x):=2\sum_{\substack{0\le i<j<m\\j-i\le7}}
                 K_f(x_j-x_i)^2
 \ge a_8(m-7)-B_8S.
 \label{eq:am-band-floor}
\end{equation}
For $m\le7$ the right side is nonpositive, so this formula also holds
for short configurations, with zero span for an empty configuration.

On a $d$-separated set the Gram matrix $U=(K_f(x_i-x_j))$ has norm
below two by \eqref{eq:am-separated-norm}. Since its diagonal is one,
\begin{equation}
 \operatorname{tr}\phi_2(U)-m
  =\operatorname{tr}(U-I)^2
  =\sum_{i\ne j}|U_{ij}|^2
  \ge a_8m-B_8S-7a_8.
 \label{eq:am-separated-defect}
\end{equation}
The entire off-diagonal energy is available on this block; selecting
the seven-band energy in \eqref{eq:am-band-floor} only makes a lower
bound smaller.

For arbitrary distinct real points, choose a maximal collection of
disjoint pairs at distance less than $d$. If two unmatched points were
closer than $d$, they could be added as another pair; thus all unmatched
points form one separated remainder. A two-point Gram block has
eigenvalues $1\pm|K_f(t)|$, both in $[0,2]$, and therefore contributes
the clipped excess $2K_f(t)^2$. By \eqref{eq:am-close-pair}, this
is at least $1/32>2a_8$.

Convex spectral pinching permits these pair blocks and the single
remainder block to be treated together. More explicitly, for any
principal-block partition, choose an orthonormal eigenbasis in each
block. Jensen's scalar inequality for the spectral measure of the
full matrix, summed over the resulting orthonormal basis, gives
$\operatorname{tr}\phi_2(U)\ge\sum_B\operatorname{tr}\phi_2(U_B)$.
The remainder has span no larger than the original span, and it alone
pays the pressure and endpoint term. We obtain
\begin{equation}
 \operatorname{tr}\phi_2(U)-n
 \ge a_8n-B_8S-7a_8
 \label{eq:am-global-affine}
\end{equation}
for every finite configuration of $n$ distinct real points.

To use this inequality with smooth trace tests, set
$\delta_\varepsilon=\|f_\varepsilon-f\|_1$. Uniformly for real $t$,
$|K_{f_\varepsilon}(t)-K_f(t)|\le\delta_\varepsilon$, and both kernels
have absolute value at most one. Each squared entry thus changes by
at most $2\delta_\varepsilon$. The ordered seven-band has at most
$14m$ entries, so its loss is at most $28\delta_\varepsilon m$.
A two-point block loses at most $4\delta_\varepsilon$.
Together with the uniform separated norm bound proved above, the
same matching argument gives
\begin{equation}
 \operatorname{tr}\phi_2(U_\varepsilon)-n
 \ge(a_8-28\delta_\varepsilon)n-B_8S-7a_8,
 \qquad\delta_\varepsilon\longrightarrow0.
 \label{eq:am-smoothed-affine}
\end{equation}
This error is proportional to the number of points. The proof never
sums uniform approximation errors over all $n^2$ matrix entries.

\subsection{The zero-counting proportion}

Apply \eqref{eq:am-smoothed-affine} to the simple critical-line part
of the signed zero operator. The two-level trace inequality with
clipping and comparison level both equal to two gives
\[
 \operatorname{tr}\mathcal A^2
 \ge2(N-n)+\operatorname{tr}\phi_2(U_\varepsilon).
\]
The scaled span is at most $N+o(N)$. Keep the smoothing fixed while
height tends to infinity, use the trace energy formula, and then let
$\varepsilon$ tend to zero. Since $a_8<1$, this yields
\begin{equation}
 \liminf_{T\to\infty}\frac{N_{\rm s0}(T)}{N(T)}
 \ge\frac{H-B_8}{1-a_8}
 \ge q:=\frac{66812491}{99194876}
 =0.673547805029767868\ldots.
 \label{eq:am-proportion}
\end{equation}
The exact rational step is
\[
 (1-a_8)\frac{66812491}{99194876}=H_0-B_8.
\]
All displayed constants satisfy the needed guards
$0<a_8<1/64$, $B_8\ge0$, and $\Lambda<2$.
In particular the simple critical-line proportion is at least
$67.35478050\%$, with that terminating percentage rounded down.


\section{What the energy says about the different counts}
\label{sec:plus-counts}

Simplicity, critical-line location, and multiplicity give different ways
to spend the same energy budget. We write $N_{\rm s0}(T)$ for simple
critical-line zeros, $N_{\rm s}(T)$ for simple zeros anywhere,
$N_0(T)$ for critical-line zeros with multiplicity, and $N_{\rm d}(T)$
for distinct complex zero locations. The count $N_{\rm s\cup0}(T)$
includes the multiplicity mass of zeros that are simple or on the line,
with the overlap counted once. Finally, $M_{\le r}(T)$ is the mass
at locations of multiplicity at most $r$, and $D_{\le r}(T)$ is the
number of those locations. Distinct locations with the same ordinate
remain distinct. We suppress $T$ inside the finite counting arguments.

The resulting lower limiting proportions are collected in
Table~\ref{tab:plus-counts}. Each complete deduction uses its own
profile and that profile's energy. The strengthened simple-block
certificate supplies the first and fourth rows. The labelled local
inequality supplies the second, and its elementary counting consequence
supplies the third. The remaining rows use the third local profile
with the multiplicity accounting developed in
\cite{Lamzouri,ChandruP6,Knausgard2026}.

\begin{table}[htbp]
\centering\small
\caption{Lower bounds for the liminf of each displayed ratio as
$T\to\infty$. Percentages are rounded down to eight decimal places.}
\label{tab:plus-counts}
\begin{tabular}{lrl}
\toprule
Ratio & Lower percentage & Input to the deduction\\
\midrule
$N_{\rm s0}/N$ & $67.35478050$ & Strengthened simple-block floor\\
$N_{\rm d}/N$ & $83.90046180$ & Labelled local floor\\
$N_{\rm s}/N$ & $67.80092360$ & Distinct-count identity\\
$(N_{\rm s}+N_0)/(2N)$ & $83.67739025$ & Retained excess\\
$N_{\rm s\cup0}/N$ & $88.93165002$ & Third-profile tangent\\
$M_{\le2}/N$ & $83.79775462$ & Third-profile mass costs\\
$D_{\le2}/N_{\rm d}$ & $93.94524765$ & Third-profile location costs\\
\bottomrule
\end{tabular}
\end{table}

The strengthened certificate has
\begin{equation}
 \widehat a=\frac{805124}{10^8},\qquad
 \widehat b=\frac{404350}{10^8},\qquad
 R_*\le R_*^+:=2-\frac{67216841}{10^8}.
 \label{eq:plus-main-count-data}
\end{equation}
Consequently its simple-critical-line proportion is
\begin{equation}
 q_*:=\frac{2-R_*^+-\widehat b}{1-\widehat a}
 =\frac{66812491}{99194876}.
 \label{eq:plus-main-count}
\end{equation}
The coefficient $\widehat a$ is positive. This matters for the
average consequence, because the finite argument preserves more
information than its final numerical bound.

\subsection{The excess and the average count}

Let $n=N_{\rm s0}$, and let $b_{\rm exc}$ count one orbit for each
multiple critical-line location and one for each reflected off-line
pair. Its excess mass is
\[
 e=N-n-2b_{\rm exc}\ge0.
\]
A multiple real location of multiplicity $j$ contributes $j-2$ to
$e$, while a reflected pair of multiplicity $j$ at each member
contributes $2(j-1)$. Before the positive-index count is replaced
by $(N-n)/2$, the clipping inequality and a simple-block certificate
with coefficients $a,b$ give
\begin{equation}
 \operatorname{tr}\mathcal A^2+bS+o(N)+O(1)
 \ge 2N-(1-a)n+2e.
 \label{eq:plus-retained-excess}
\end{equation}
Here the smoothing error is removed after the height limit, as in
the preceding construction.

Let $B_{\rm off}$ be the mass of multiple off-line zeros.
The contribution of a pair with multiplicity $j\ge2$ is $2j$
to $B_{\rm off}$ and $2(j-1)$ to $e$, so $B_{\rm off}\le2e$.
Inclusion--exclusion for the two multiplicity counts is the exact identity
\[
 N_{\rm s}+N_0=N+n-B_{\rm off}.
\]
Combining it with \eqref{eq:plus-retained-excess} yields
\[
 N_{\rm s}+N_0
 \ge 3N-\operatorname{tr}\mathcal A^2-bS+an-o(N)-O(1).
\]
For $a\ge0$, insert the simple-critical lower bound
$q=(2-R^+-b)/(1-a)$ and the trace and span limits. The identity
$3-R^+-b+aq=1+q$ then gives
\begin{equation}
 \liminf_{T\to\infty}\frac{N_{\rm s}(T)+N_0(T)}{2N(T)}
 \ge\frac{1+q}{2}.
 \label{eq:plus-average-transfer}
\end{equation}
At \eqref{eq:plus-main-count} the right side is
$166007367/198389752$. This is a lower bound for one average;
it does not separately bound either component by that number.
The retained-excess deduction is part of the earlier counting
framework \cite{ChandruP6,Knausgard2026}.

\subsection{Profile energies and the labelled distinct count}

For a finite cosine profile
$f(u)=1+\sum_{j=1}^D a_j\cos(2\pi ju)$ on $[-1/2,1/2]$,
direct integration gives
\begin{equation}
 R(f)=\frac43+\frac12\sum_{j=1}^D a_j^2+
 \frac1{\pi^2}\sum_{j=1}^D
       \frac{(-1)^ja_j-a_j^2/4}{j^2}.
 \label{eq:plus-cosine-energy}
\end{equation}
The first and third profiles of \cite{Knausgard2026} have coefficient
lists, in units of $10^{-6}$,
\[
 (193305,-23105,29166,13952,-21009,41052),\qquad
 (200566,19,-9358,9283,-6746,8170).
\]
Rational enclosures in \eqref{eq:plus-cosine-energy} give
\begin{equation}
 R_1\le R_1^+:=\frac{666352479141}{500000000000},\qquad
 R_3\le R_3^+:=\frac{83278977073}{62500000000}.
 \label{eq:plus-other-energy-caps}
\end{equation}
These evaluate the existing profiles more closely; their local
inequalities are the cited computer-assisted inputs.

For the distinct count the first profile is used with
\[
 c=\frac{67}{20},\qquad
 x_1=\frac{54350}{10^6},\quad
 x_2=\frac{69433}{10^6},\quad b_1=\frac{37314}{10^6}.
\]
If $U$ is the Gram matrix, formed with the smoothed kernel, of real
locations carrying labels $d_i\in\{1,2\}$, set $G=D^{1/2}UD^{1/2}$, where
$D=\operatorname{diag}(d_i)$. The labelled input is
\[
 \operatorname{tr}\phi_c(G)-\operatorname{tr}D^2
 \ge x_1n_1+x_2n_2-b_1S-\omega(\varepsilon)(n_1+n_2)-O(1).
\]
Here $n_1,n_2$ count simple and double real locations, and
$\omega(\varepsilon)\to0$; the height limit precedes removal of the
smoothing. Retain all real locations in the positive part, pinch the
higher multiplicities into scalar blocks, and leave reflected off-line
pairs in the remainder.
Write $h$ for the number of higher-multiplicity real locations and
$k$ for the number of off-line pairs, so
$N_{\rm d}=n_1+n_2+h+2k$.

Put
\[
 \begin{aligned}
 \mu&=2-2x_1+x_2,&\nu&=2c-3+x_1-x_2,\\
 r_h&=6c-7-c^2+x_1-2x_2,&
 r_k&=4c-2-c^2-2x_1.
 \end{aligned}
\]
Their values are
\[
 (\mu,\nu,r_h,r_k)=
 \left(\frac{1960733}{10^6},\frac{3684917}{10^6},
       \frac{224123}{125000},\frac{43}{625}\right),
\]
all positive. The energy debit is at most
$(\nu+\mu)n_1+(2\nu+\mu)n_2+c^2(h+k)$.
It is bounded by $\nu N+\mu N_{\rm d}$: a real location of
multiplicity $j\ge3$ leaves the nonnegative difference
$j\nu+\mu-c^2\ge r_h$, and an off-line pair of multiplicity
$j\ge1$ leaves $2j\nu+2\mu-c^2\ge r_k$.
The energy budget therefore gives
\begin{equation}
 \liminf\frac{N_{\rm d}}N
 \ge\frac{3-R_1^++x_2-x_1-b_1}{2-2x_1+x_2}
 =\frac{822532020859}{980366500000}.
 \label{eq:plus-distinct}
\end{equation}
At each location, multiplicity one contributes one to
$2N_{\rm d}-N$, multiplicity two contributes zero, and higher
multiplicities contribute negatively. Thus
$N_{\rm s}\ge2N_{\rm d}-N$, and
\begin{equation}
 \liminf\frac{N_{\rm s}}N
 \ge\frac{332348770859}{490183250000}.
 \label{eq:plus-simple-anywhere}
\end{equation}
This last count includes zeros off the critical line.

\subsection{The union and the full fixed-cutoff hierarchy}

For the remaining counts use the third profile's simple-block floor,
with
\begin{equation}
 a=\frac{54525}{10^6},\qquad b=\frac{37272}{10^6},\qquad
 \tau=R_3^++b=\frac{85608477073}{62500000000},\qquad
 \beta=\tau-1-a=\frac{19700664573}{62500000000}.
 \label{eq:plus-third-budget}
\end{equation}
Define
\begin{equation}
 c_*=2+\sqrt{2-2a},\qquad
 Z=\frac{c_*^2}{2}=3-a+2\sqrt{2-2a},\qquad E=\frac\beta Z.
 \label{eq:plus-excess-budget}
\end{equation}
In the finite clipping inequality use equal levels $c_*$. The
simple-block floor remains valid because $\phi_c(s)$ increases with
$c$ for $s\ge0$, and $c_*>2$. In terms of $u=N-n$ and
$e=u-2b_{\rm exc}$, the resulting energy lower bound is
\[
 (1+a)N+
 \left(2c_*-1-a-\frac{c_*^2}{2}\right)u
 +\frac{c_*^2}{2}e-bS-o(N)-O(1).
\]
The coefficient of $u$ vanishes exactly. Consequently
\begin{equation}
 \limsup_{T\to\infty}\frac{e(T)}{N(T)}\le E.
 \label{eq:plus-excess-limit}
\end{equation}
As $N-N_{\rm s\cup0}\le2e$, this proves
\begin{equation}
 \liminf\frac{N_{\rm s\cup0}}N\ge1-2E
 =1-\frac{2(\tau-1-a)}{3-a+2\sqrt{2-2a}}.
 \label{eq:plus-union}
\end{equation}
Equivalently, $c_*$ is the largest permissible tangent level for
which simple off-line pairs still cost at least $1+a$ per unit mass;
the cost of multiple off-line pairs increases up to that level.

The same excess controls every fixed higher cutoff. A real location
of multiplicity $j\ge r+1$ has mass-to-excess ratio
$j/(j-2)\le(r+1)/(r-1)$. A nonreal pair has the smaller ratio
$j/(j-1)$. Hence, for every fixed integer $r\ge3$,
\begin{equation}
 \liminf\frac{M_{\le r}}N
 \ge1-\frac{r+1}{r-1}E.
 \label{eq:plus-higher-mass}
\end{equation}
This is the retained multiplicity-budget mechanism of
\cite{Waddle,ChandruP6}, evaluated with \eqref{eq:plus-third-budget}.
It uses the excess estimate, rather than the union bound alone.

For the corresponding location proportions, and for the stronger
cutoff-two result, it is useful to give the common accounting once.
Let $x_j$ count real locations of multiplicity $j$, and $y_j$ count
reflected pairs of multiplicity $j$ at each member. Then
\[
 N=\sum_jjx_j+2\sum_jjy_j,\qquad
 N_{\rm d}=\sum_jx_j+2\sum_jy_j.
\]
Put $h_c(s)=2cs-\phi_c(s)$. Retain all real locations in the positive
part and pinch their Gram matrix into the simple block and scalar
multiple-location blocks. For every $c\ge2$, the finite inequality is
\begin{equation}
 \operatorname{tr}\mathcal A^2+bS+o(N)+O(1)
 \ge 2cN-\sum_jh_c(j)x_j-c^2\sum_jy_j+ax_1.
 \label{eq:plus-multiplicity-master}
\end{equation}
All approximation errors can first be included in a fixed upper
budget $\tau' N$, with $\tau'>\tau$. This also justifies the
subsequent division by the distinct count.

Set $\alpha=2c-\tau'$ and $\beta'=\tau'-1-a$. Moving the budget
to the right gives
\[
 \sum_j\lambda_jx_j+\sum_j\omega_jy_j\le0,\qquad
 \lambda_j=\alpha j-h_c(j)+a\mathbf1_{j=1},\quad
 \omega_j=2\alpha j-c^2.
\]
If each low-multiplicity location costs at least $-\beta'$ and each
high-multiplicity location at least $\gamma_r-\beta'$, then
\[
 0\ge-\beta'N_{\rm d}+\gamma_r(N_{\rm d}-D_{\le r}),
 \qquad
 \frac{D_{\le r}}{N_{\rm d}}\ge1-\frac{\beta'}{\gamma_r}.
\]
A reflected pair carries two location costs.

For $r\ge3$, choose $c=c_*\in(3,4)$. At the limiting budget,
$\alpha=Z-\beta$ and
$2c_*-1-c_*^2/2-a=0$. The real-location costs plus $\beta$ at
$j=1,2,3$ are respectively
$0$, $3-\tau-a$, and $8-2\tau-a$.
For $j\ge4$ the expression is $\alpha j+\beta-c_*^2$.
The required guards are
\[
 \alpha>0,\quad\beta>0,\quad
 3-\tau-a>0,\quad8-2\tau-a>0,\quad
 4\alpha+\beta-c_*^2=2Z-3\beta>0.
\]
They hold here: $a<1/2$, $\tau<7/5$, $0<\beta<1/3$, and
$Z>9/2$ suffice. The first high real location has cost plus $\beta$
equal to
\[
 \gamma_r=(r+1)\alpha+\beta-c_*^2
 =Z(r-1-rE)>0.
\]
For pairs, the lowest cost plus $2\beta$ is zero at $j=1$,
and the high-pair cost at $j=r+1$ exceeds $2\gamma_r$ by $c_*^2$;
all later costs increase. The strict guards persist for
$\tau'>\tau$ sufficiently close to $\tau$, while the zero guard
for the simple pair is independent of $\tau'$. Letting
$\tau'\downarrow\tau$ gives
\begin{equation}
 \liminf\frac{D_{\le r}}{N_{\rm d}}
 \ge1-\frac{E}{r-1-rE}\qquad(r\ge3\text{ fixed}).
 \label{eq:plus-higher-locations}
\end{equation}
These fixed-cutoff consequences are most useful at small cutoffs.
For sufficiently large cutoffs, the unconditional exponential
multiplicity estimates of \cite{STT} give stronger mass tails;
the right side of \eqref{eq:plus-higher-mass} instead tends to $1-E$.

\subsection{Multiplicity at most two and the sharp scalar population}

At $c=3$ the real-location debits in
\eqref{eq:plus-multiplicity-master} are $5-a$, $8$, and $9$ for
$j=1$, $j=2$, and $j\ge3$. Every off-line pair has debit nine.
Since $0\le a\le1/2$, the debit per unit mass at multiplicity at
most two is at most $5-a$; for the remaining mass it is at most
three. Therefore
\begin{equation}
 \liminf\frac{M_{\le2}}N
 \ge\frac{3-\tau}{2-a}
 =\frac{101891522927}{121592187500}.
 \label{eq:plus-mass-two}
\end{equation}

For location costs take $\alpha=6-\tau'$,
$\beta'=\tau'-1-a$, and $\gamma_2=8-2\tau'-a$.
The sufficient guards are
\[
 \beta'\ge0,\quad\alpha>0,\quad
 3-\tau'-a\ge0,\quad1-2a\ge0,\quad\gamma_2>0.
\]
The real costs satisfy $\lambda_1=-\beta'$,
$\lambda_2+\beta'=3-\tau'-a\ge0$, and
$\lambda_j+\beta'-\gamma_2=\alpha(j-3)\ge0$ for $j\ge3$.
For pairs, $\omega_1+2\beta'=1-2a\ge0$ and
$\omega_3+2\beta'-2\gamma_2=9$, after which the costs increase.
All guards are strict at \eqref{eq:plus-third-budget}. The location
argument followed by $\tau'\downarrow\tau$ gives
\begin{equation}
 \liminf\frac{D_{\le2}}{N_{\rm d}}
 \ge1-\frac{\tau-1-a}{8-2\tau-a}
 =\frac{101891522927}{108458411118}.
 \label{eq:plus-locations-two}
\end{equation}

There is a direct reason that these two expressions fit together.
Within the scalar orbit accounting of
\eqref{eq:plus-multiplicity-master}, they are attained simultaneously.
Assign a fraction
$\mu=(3-\tau)/(2-a)$ of the mass to simple real locations, and
the remaining fraction to triple real locations. The energy charge
is
\[
 (1+a)\mu+3(1-\mu)=\tau.
\]
The fraction of low-multiplicity locations is then
\[
 \frac{\mu}{\mu+(1-\mu)/3}
 =\frac{3\mu}{1+2\mu}
 =1-\frac{\tau-1-a}{8-2\tau-a}.
\]
This population satisfies every scalar row obtained from
\eqref{eq:plus-multiplicity-master} by changing $c\ge2$:
$\phi_c(1)=1$ and $\phi_c(3)\le9$, with equality in the latter
for $c\ge3$. Its proportions are rational, so a common denominator
produces integral location counts. Thus changing the clipping level
alone cannot strengthen these two conclusions within that accounting.
The population certifies sharpness of the reduced scalar constraints;
it assigns no Gram interactions between its locations. Additional
geometric or arithmetic information can impose further constraints.

The rational values in this section and the union radical in
\eqref{eq:plus-union} determine all entries of
Table~\ref{tab:plus-counts}. In particular, the larger simple-anywhere
number, the average number, and the distinct-location number each
retain their own numerator and denominator throughout the argument.


\section{Local slack and continuous payment for long spans}
\label{sec:rigidity}

\subsection{The fixed local input and its retained slack}

The main numerical bound is now established. We turn to a different
question: what quantitative information is retained by local slack and
by pair energy beyond the lags used in a local floor? An explicit
application uses the second profile of
\cite[Section 5]{Knausgard2026}:
\begin{equation}
 \begin{aligned}
 f(u)&=1+\sum_{j=1}^{10}a_j\cos(2\pi ju),\\
 10^6(a_1,\ldots,a_{10})&=
 (122290,-45012,16241,1887,2711,3731,\\
 &\hspace{1.2em}-5307,2982,-5153,5358).
 \end{aligned}
 \label{eq:profile-two}
\end{equation}
The coefficient absolute sum is $0.210672$, so $f>0$ and $|K_f|\le1$.
Throughout this section $K=K_f$, $\theta=17/20$, and
\begin{equation}
 a=\frac{8722}{10^6},\quad \alpha=\frac{627}{10^6},\quad
 b=7\alpha=\frac{4389}{10^6},\quad
 d_0=\frac{10016}{10^6},\quad
 B_0=7a+\frac{3882}{10^6}.
 \label{eq:base-constants}
\end{equation}
The public eight-point certificate, with weights and bounded corrections
listed in Appendix~\ref{app:local-data}, says that for every seven-gap
frame with gaps at least $\theta$,
\begin{equation}
 F_i:=\alpha(y_{i+7}-y_i)+
 \sum_{0\le r<s\le7}W_{rs}K(y_{i+s}-y_{i+r})^2
 +\sum_{r=0}^6B_r(g_{i+r})\ge a.
 \label{eq:local-floor}
\end{equation}
Its validity is the imported computer-assisted input
\cite[Proposition 5.3]{Knausgard2026}. Define
$\ell_i=F_i-a\ge0$ and $H=\sum_i\ell_i$ over all complete frames.
With ordered short-band energy
\[
 E_s=2\sum_{h=1}^7\sum_iK(y_{i+h}-y_i)^2,
\]
the pair-capacity sums are at most two. The spread charge sums to at
most $7\alpha S$; the bounded corrections telescope with loss at most
$3882/10^6$. Retaining, rather than dropping, the slack gives
\begin{equation}
 E_s\ge am-bS-B_0+H,\qquad
 \ell_i\ge\alpha(y_{i+7}-y_i)-d_0.
 \label{eq:retained-slack}
\end{equation}
The second inequality uses nonnegative pair weights and
$\sum_r\min B_r=-1294/10^6$. For fewer than eight points take
$H=0$; the first inequality still holds because $B_0\ge7a$.

\subsection{A general phase comparison}

The mechanism is more general than these data. Let $K$ be real-valued
on the real axis. Suppose frames of $h$
gaps have nonnegative slack satisfying
$\ell_i\ge\alpha(y_{i+h}-y_i)-d_0$, with $\alpha>0$.
Assume gaps are at least $\theta>0$ and
\begin{equation}
 |K(t)|\ge\frac{h_0|\sin\pi t|}{\pi t}\quad(t\ge t_0),
 \qquad h_0>0.
 \label{eq:tail-hypothesis}
\end{equation}
Choose finitely many integer lags $l\ge h$ with $l\theta\ge t_0$
and nonnegative weights $w_l$ satisfying
\begin{equation}
 |e^{ix}-1|^2\le\sum_lw_l|e^{ilx}-1|^2\quad(x\in\R).
 \label{eq:phase-comparison}
\end{equation}
Put $q_l=\lceil l/h\rceil$ and choose positive numbers
\begin{equation}
 R_l\ge q_ld_0/\alpha-(hq_l-l)\theta.
 \label{eq:cover-radius}
\end{equation}
Writing $T=2\sum_l\sum_i K(y_{i+l}-y_i)^2$, one obtains
\begin{equation}
 \sum_i\dist(g_i,\Z)^2\le M_ET+M_HH+B_1,
 \label{eq:general-rigidity}
\end{equation}
where
\begin{equation}
 M_E=\max_l\frac{w_l\pi^2R_l^2}{8h_0^2},\qquad
 M_H=\sum_l\frac{w_lq_l}{2\alpha R_l},\qquad
 B_1=\frac18\sum_lw_ll.
 \label{eq:prices}
\end{equation}
Here and below empty sums are zero.

To prove this, set $c(t)=|e^{2\pi it}-1|^2\le4$.
For every $R>0$ and every $t\ge t_0$, \eqref{eq:tail-hypothesis} gives
\begin{equation}
 c(t)\le\frac{4\pi^2R^2}{h_0^2}K(t)^2+
 \frac8R(t-R)_+.
 \label{eq:continuous-payment}
\end{equation}
For $t\le R$ the first term suffices. For $t>R$ it pays at least
$c(t)R^2/t^2$; the remainder is at most
$4(1-R^2/t^2)\le8(t-R)/R$.

A lag-$l$ edge is covered by the $q=q_l$ complete frame starts
$0,h,\ldots,(q-2)h,l-h$ relative to its anchor, with a single frame
when $l=h$. There are exactly $hq-l$ repeated gap incidences.
Since each gap is at least $\theta$, the sum $L_i$ of the cover slacks
satisfies
\[
 L_i\ge0,\qquad
 L_i\ge\alpha t+\alpha(hq-l)\theta-qd_0
 \ge\alpha(t-R_l).
\]
Each fixed frame occurs in at most $q_l$ translated covers for that lag.
Summing \eqref{eq:continuous-payment} over its anchors consequently gives
\begin{equation}
 \sum_i c(y_{i+l}-y_i)
 \le\frac{4\pi^2R_l^2}{h_0^2}\sum_iK(y_{i+l}-y_i)^2
 +\frac{8q_l}{\alpha R_l}H.
 \label{eq:covered-chords}
\end{equation}
All frames lie inside their edge, and every physical span is included.

Now extend $z_i=e^{2\pi iy_i}$ by zero outside the finite configuration.
Parseval applies \eqref{eq:phase-comparison} to its discrete differences.
For lag $l$ the zero extension adds $2\min(l,m)\le2l$ unit boundary
terms. Dropping the favorable nearest-neighbor boundary gives
\[
 \sum_i|z_{i+1}-z_i|^2
 \le\sum_lw_l\sum_i|z_{i+l}-z_i|^2+2\sum_lw_ll.
\]
Finally $\dist(g,\Z)^2\le|e^{2\pi ig}-1|^2/16$.
Insert \eqref{eq:covered-chords} and retain the factor two in the
ordered energy $T$ to obtain \eqref{eq:general-rigidity}.

There is a simple uniform family of phase comparisons. On $[r,2r]$
take endpoint weights $2/r$ and interior weights $4/r$.
For the unitary shift $V$,
$V-I=(V^{l+1}-I)-V(V^l-I)$; the squared triangle inequality averaged
over $r\le l<2r$ proves \eqref{eq:phase-comparison}. The weight sum
is four and its first moment is $6r$. If $d_0>0$, the choice
$R_l=q_ld_0/\alpha$ gives
\begin{equation}
 M_H=2/d_0,\qquad
 M_E\le\frac{\pi^2d_0^2\lceil2r/h\rceil^2}
 {2h_0^2\alpha^2r},\qquad B_1=3r/4.
 \label{eq:linear-horizon}
\end{equation}
Thus the energy price grows at most linearly in $r$, while the slack
price is independent of $r$. If the local floor spends every lag through
$h$, take $r>h$ before adding $T$ to that floor.

For degree-$D$ cosine profiles with
$\sum|a_j|\le A_0$ and $f(1/2)\ge v_0>0$, the rational sinc factor
satisfies
\[
 V(t^2)=f(1/2)+\sum_{j=1}^D(-1)^ja_j\frac{j^2}{t^2-j^2},\qquad
 |V(t^2)-f(1/2)|\le2A_0D^2/t^2\quad(t^2\ge2D^2).
\]
Hence $h_0=v_0/2$ is available beyond
$D\max(\sqrt2,2\sqrt{A_0/v_0})$. More precisely, if
$h,\theta,\alpha,d_0,A_0,v_0$ are fixed, with
$\theta,\alpha,d_0,v_0>0$, a family satisfying the stated frame
hypothesis permits $r>h$ with $r\theta$ beyond that threshold and
$r=O(D)$ for $D\ge1$. Formula~\eqref{eq:linear-horizon} then gives
an $O(D)$ energy price and a constant slack price.
Each profile and its finite lag set are fixed before the height limit.

\subsection{An exact finite instance}

For \eqref{eq:profile-two}, write $c_j=(-1)^ja_j$ and
$z_0=(51/5)^2$. Its sinc factor is
\[
 K(t)=\frac{\sin\pi t}{\pi t}V(t^2),\qquad
 V(z)=1+\sum_{j=1}^{10}c_j\frac{z}{z-j^2}.
\]
For negative $c_j$ use the value at $z_0$; for positive $c_j$ use
the limiting value $c_j$. This proves, for $t\ge51/5$,
\begin{equation}
 V(t^2)\ge\frac{3123731672835001}{3750959784000000}
 >\frac{104}{125}=:h_0.
 \label{eq:tail-factor}
\end{equation}
Take
\begin{equation}
 \mathcal L=\{12,\ldots,49\},\quad q_l=\lceil l/7\rceil,\quad
 R_l=\frac{201q_l+17l}{20},\quad
 w_l=\frac{196}{R_l^2}.
 \label{eq:lag-data}
\end{equation}
The identity $R_l=16q_l-(7q_l-l)\theta$ and
$16>d_0/\alpha=10016/627$ verify \eqref{eq:cover-radius}.
The phase comparison is certified by the polynomial
\begin{equation}
 P(t)=\frac{\sum_{l=12}^{49}w_l(1-T_l(t))-(1-t)}{1-t},
 \label{eq:phase-polynomial}
\end{equation}
where $T_l$ is the Chebyshev polynomial. Its numerator vanishes at one,
so $P$ has degree 48. All 49 Bernstein coefficients on each of the
following intervals exceed the displayed rational bound:
\begin{table}[htbp]
\centering
\caption{Certified lower bounds for every Bernstein coefficient of the
phase polynomial on six intervals.}
\label{tab:phase-bounds}
\begin{tabular}{cc}
\toprule Interval & Lower bound\\\midrule
$[-1,-3/4]$ & $164977/50000000$\\
$[-3/4,-1/2]$ & $255867387/1000000000$\\
$[-1/2,0]$ & $370473121/1000000000$\\
$[0,1/2]$ & $207778923/200000000$\\
$[1/2,3/4]$ & $743826529/250000000$\\
$[3/4,1]$ & $1129786507/125000000$\\\bottomrule
\end{tabular}
\end{table}
The nonnegative Bernstein basis proves $P\ge0$ on $[-1,1]$ and
therefore \eqref{eq:phase-comparison} on the whole real line.
Appendix~\ref{app:arithmetic} specifies the exact basis conversion.
The prices in \eqref{eq:prices}, using $\pi<22/7$, satisfy
\begin{equation}
 M_E<\frac{1890625}{5408}<350,\qquad
 M_H<112,\qquad B_1<6.
 \label{eq:frozen-prices}
\end{equation}
With $J=H+T$, \eqref{eq:general-rigidity} gives
\eqref{eq:intro-rigidity}. The unused lags $12$--$49$ are disjoint
from the paid lags $1$--$7$, so $D=E_s+T$ also satisfies
\begin{equation}
 D\ge am-bS-B_0+J.
 \label{eq:retained-J}
\end{equation}

\subsection{A signed pressure row for the retained phase resource}
\label{sec:bootstrap}


For every $g\ge17/20$ we prove
\begin{equation}
 2K(g)^2+pg+C\dist(g,\Z)^2\ge c,
 \qquad p=\frac{17}{4000},\quad C=\frac{19}{4},\quad
 c=\frac{1891}{200000}.
 \label{eq:pressure}
\end{equation}
The profile bound $f\le1.210672$ gives
$|K''|\le\pi^2(1.210672)/3<4$.
The sinc expansion
\[
 K(t)=\sinc t+\tfrac12\sum_{j=1}^{10}a_j
 \{\sinc(t-j)+\sinc(t+j)\},\qquad
 \sinc t=\frac{\sin\pi t}{\pi t},
\]
gives the exact values
\begin{center}
\begin{tabular}{ccc}
\toprule $k$ & $K(k)$ & $K'(k)$\\\midrule
$1$ & $12229/200000$ & $-260321359/264000000$\\
$2$ & $-11253/500000$ & $2577793683/6160000000$\\\bottomrule
\end{tabular}
\end{center}
Write $A_k=|K(k)|$, $B_k=|K'(k)|$ and take
$r_1=1/20$, $r_2=1/25$. At both integers the signed derivative of
$|K|$ is negative. Taylor's formula gives, for $d\ge0$,
\[
 \operatorname{sign}(K(k))K(k+d)\ge A_k-B_kd-2d^2,
 \qquad
 \operatorname{sign}(K(k))K(k-d)\ge A_k+B_kd-2d^2.
\]
On $0\le d\le r_k$ the first right side is positive. For $g=k+d$
the left side of \eqref{eq:pressure} minus $c$ is at least
\[
 q_0+q_1d+q_2d^2+8B_kd^3+8d^4,
 \quad q_0=2A_k^2+pk-c,\quad
 q_1=-4A_kB_k+p,\quad q_2=2B_k^2-8A_k+C.
\]
Here $q_2>0$ and the exact positive margins $4q_0q_2-q_1^2$ are
\begin{equation}
 \frac{53039685134812231}{193600000000000000000},\qquad
 \frac{23907572162979771241}{948640000000000000000000}
 \label{eq:discriminants}
\end{equation}
for $k=1,2$. Completing the quadratic square proves the right-hand
inner intervals. On the left,
$|K(k-d)|\ge A_k+(B_k-2r_k)d$ and
\[
 2K(k-d)^2+p(k-d)+Cd^2-c
 \ge q_0+\{4A_k(B_k-2r_k)-p\}d.
\]
The two linear coefficients are respectively
\[
 \frac{2804524299211}{13200000000000},\qquad
 \frac{1835448537709}{70000000000000},
\]
both positive.

For $r_k\le d\le1/2$, discard the kernel term. The lower expression
$p(k-d)+Cd^2-c$ increases with $d$ and at the two cutoffs equals
$2583/400000$ and $259/40000$. If a nearest integer is at least
three, then $g\ge5/2$ and $pg>c$. These cases cover every permitted
gap, including half-integer ties, and prove \eqref{eq:pressure}
without sampled-grid acceptance.

\subsection{Cancellation of a single retained resource}

The following algebra explains the role of phase rigidity. Suppose
$D\ge am-bS-B_0+J$, $J\ge0$, and
$\sum\dist(g_i,\Z)^2\le MJ+B_1$, while
$D\ge2\sum_iK(g_i)^2$. Summing a row of the form
\eqref{eq:pressure} gives the alternative bound
\[
 D\ge cm-pS-c-CB_1-MCJ.
\]
Taking convex weights $MC/(1+MC)$ and $1/(1+MC)$ cancels $J$.
The resulting coefficients are
\begin{equation}
 a'=a+\frac{c-a}{1+MC},\qquad
 b'=b+\frac{p-b}{1+MC},\qquad
 B'=\frac{MCB_0+c+CB_1}{1+MC}.
 \label{eq:bootstrap}
\end{equation}
This combines alternative bounds on the same energy; it does not add
overlapping pair resources.

For $M=350$, $B_1=6$ and \eqref{eq:pressure},
$1+MC=3327/2$ and
\begin{equation}
 D\ge a_\star m-b_\star S-B_\star,\qquad
 B_\star=\frac{27293111}{332700000},
 \label{eq:final-floor}
\end{equation}
where
\[
 a_\star=\frac{725489}{83175000},\qquad
 b_\star=\frac{584077}{133080000}.
\]
There are 45 retained index bands. Under
$\sup|K_\varepsilon-K|\le\eta\le1$,
$|K_\varepsilon^2-K^2|\le3\eta$.
Keep $J$ defined on the original kernel. The first floor loses at
most $270\eta m$; the adjacent-gap row loses at most $6\eta m$.
Their same convex combination therefore gives
\begin{equation}
 D(K_\varepsilon)\ge
 \left(a_\star-\frac{299254}{1109}\eta\right)m
 -b_\star S-B_\star.
 \label{eq:final-smoothing}
\end{equation}

The imported majorant for this profile has
$\Lambda=1.78942\ldots<2$, and the near-pair bound is
$\beta=13707941/2^{28}>a_\star$ \cite[Section 5]{Knausgard2026}.
Choose $\eta<(\beta-a_\star)/3$. Each removed pair pays more than
$2a_\star$, the separated block is unclipped, and $b_\star>0$.
The matching argument of Section~\ref{sec:am-transfer} transfers
\eqref{eq:final-smoothing} to every real configuration. The coefficients
satisfy $a_\star>a$ and $0<b_\star<b$: the same local input has acquired
a larger point reward and a smaller spread cost. This is a quantitative
use of the retained phase resource, for the profile
\eqref{eq:profile-two}; it is not an additional gain in
\eqref{eq:intro-main}, which uses a different profile.

\section{A common residual for joint affine deductions}
\label{sec:structural}

The phase comparison retains local slack and unused pair energy
on the same configuration.
We now ask what can be deduced when the retained information has
already been reduced to universal affine floors. The distinction is
between precisely specified sets of inequalities, not between
different assertions about the zeros themselves.

\subsection{The relaxation being tested}

Continue to write $I=[-1/2,1/2]$. For a real function $f\in L^2(I)$,
extended by zero, put
\begin{equation}
 \begin{split}
 A(f)&=\int_I f(s)^2\dd s,\\
 J(f)&=\iint_{I^2}|s-t|f(s)f(t)\dd s\dd t,\\
 C_h(f)&=\int_{\R}f(s)f(s+h)\dd s,\qquad R(f)=A(f)+J(f).
 \end{split}
 \label{eq:structural-forms}
\end{equation}
The variational arguments will allow signed real functions with
$\int f=1$. A profile used in a Gram matrix is nonnegative; an analytic
profile is also even and is subject to the fixed-test approximation in
Section~\ref{sec:interface}. Bounds proved on this larger space
apply in particular to the nonnegative subclass.

For a normalized nonnegative profile let
$G_f=(K_f(x_j-x_k))_{j,k=1}^n$ on a real configuration.
It is a positive semidefinite Hermitian matrix with unit diagonal,
including when $f$ is not even. For $t\ge0$ define
\begin{equation}
 D_t(G_f)=
 \tr\bigl((G_f-I)^2-(G_f-(1+t)I)_+^2\bigr).
 \label{eq:structural-defect}
\end{equation}
The scalar function in this definition is nonnegative on $[0,\infty)$,
and the subtracted square is nonnegative. Therefore
\begin{equation}
 0\le D_t(G_f)\le\tr(G_f-I)^2
 =\sum_{j\ne k}|K_f(x_j-x_k)|^2.
 \label{eq:defect-frobenius}
\end{equation}
Since $\tr G_f=n$, we also have
$D_t(G_f)=\tr\phi_{1+t}(G_f)-n$.

Let $b_{\rm exc}$ count the exceptional conjugacy orbits:
one for each multiple real location and one for each nonreal pair.
The orbit argument in Section~\ref{sec:interface} gives
$n_+(Q)\le b_{\rm exc}$ when the simple-real part is removed.
In \eqref{eq:two-level} take the two levels equal to $\Gamma/2$,
where $\Gamma\ge2$. It follows that
\begin{equation}
 \tr\mathcal A_f^2
 \ge\Gamma N-(\Gamma-1)n-\frac{\Gamma^2}{4}b_{\rm exc}
       +D_{\Gamma/2-1}(G_f).
 \label{eq:structural-inertia}
\end{equation}
This is the compensated inertia inequality of \cite{ChandruP6}.
It uses the same original multiset in every profile realization.
The exceptional-orbit count is common to the inequalities;
no rank statement for a sum of different profile operators is assumed.

Write the normalized omitted mass and excess as
\begin{equation}
 u=\frac{N-n}{N},\qquad
 e=\frac{N-n-2b_{\rm exc}}N,\qquad 0\le e\le u\le1.
 \label{eq:normalized-excess}
\end{equation}
When passing to limiting rows, take a common convergent subsequence
of these bounded ratios and use the same symbols for its limit.
This preserves one count point while the fixed-test error terms
vanish.
Nonnegativity of $e$ follows orbit by orbit: a multiple real location
of multiplicity $k$ contributes $k-2$, and a nonreal pair with
multiplicity $j$ at each member contributes $2(j-1)$.

A row consists of finitely many fixed profiles $f_i$,
weights $w_i\ge0$ with $\sum_iw_i=1$, and fixed thresholds
$\Gamma_i\ge2$. Its finite hypothesis is a genuinely joint floor
\begin{equation}
 \sum_iw_iD_{\Gamma_i/2-1}(G_{f_i})
 \ge an-dS-B,\qquad d\ge0,
 \label{eq:joint-affine-floor}
\end{equation}
valid on every real configuration of size $n$ and span $S$, with
$B$ independent of $n$. The same points occur in every matrix.
No decomposition into separately minimized floors is required.
All parameters, including the finite family and endpoint constant,
are fixed before the height limit.

Average \eqref{eq:structural-inertia} with these weights.
The trace formula and span bound of Section~\ref{sec:interface}
give the normalized necessary row
\begin{equation}
 \overline R+d\ge
 (1-u)(1+a)+\overline\Gamma u
 -\frac{u-e}{8}\sum_iw_i\Gamma_i^2,
 \quad
 \overline R=\sum_iw_iR(f_i),\quad
 \overline\Gamma=\sum_iw_i\Gamma_i.
 \label{eq:joint-master-row}
\end{equation}
The sign $d\ge0$ is needed when replacing $S/N$ by its upper limit
one. For this analytic implication each profile is a fixed admissible
smooth test, or a limit of such tests for which the finite-floor
error is uniform in the configuration and any required norm control
has been proved separately, as in Section~\ref{sec:interface}.
Continuity of matrix entries and spectral positive parts passes
such a common finite floor to its $L^2$ limit on each fixed
configuration. It does not transfer an arbitrary unsmoothed floor
back to its approximants uniformly in the number of points.
The periodic and variational arguments below are valid for all
$L^2$ profiles without this additional analytic admissibility issue.

The relaxation studied here retains any collection of rows
\eqref{eq:joint-master-row}, their common count variables, and the
orbit accounting for the desired statistic. Each row may have its
own profiles, weights, slope and thresholds. It does not automatically
retain labelled-block inequalities, square-completion slack, signed
identities between exact energies, products of operators, or further
arithmetic information. The conclusions concern precisely this
relaxation. They impose neither an upper bound on an actual
zero proportion nor a ceiling on every deduction from pair correlation.
In particular they do not apply to the labelled distinct-zero scheme
merely by changing the numerator.

\subsection{Periodic tests without a regularity assumption}

Take $p$ distinct points $x_1,\ldots,x_p$ modulo a period $L>0$.
Their full periodic energy per point, including the diagonal, is
\begin{equation}
 Q_{L,x}(f)=\frac1p\sum_{a,b=1}^p\sum_{k\in\Z}
 |K_f(x_b-x_a+kL)|^2.
 \label{eq:periodic-full-energy}
\end{equation}
The sampling identity is valid for every real $L^2$ profile.
With the Fourier convention \eqref{eq:normalization},
periodize $f(s)e^{2\pi ids}$ on an interval of length $1/L$:
\[
 P_d(s)=\sum_{\nu\in\Z}
 f(s+\nu/L)e^{2\pi id(s+\nu/L)},\qquad 0\le s<1/L.
\]
Only finitely many translates overlap, so $P_d\in L^2$.
Its normalized Fourier coefficient of index $k$ is $LK_f(d-kL)$.
Parseval therefore gives
\[
 \sum_k|K_f(d+kL)|^2
 =\frac1L\int_0^{1/L}|P_d(s)|^2\dd s.
\]
Expand the square, group by the difference of the translate indices,
and sum over $d=x_b-x_a$. This proves
\begin{equation}
 Q_{L,x}(f)=\frac1{pL}\sum_{j\in\Z}
 \left|\sum_{a=1}^p e^{2\pi ijx_a/L}\right|^2 C_{j/L}(f).
 \label{eq:periodic-autocorrelation}
\end{equation}
This is a finite sum, since $C_h(f)=0$ for $|h|\ge1$.
There is no endpoint term at $|h|=1$, where the supports intersect
only on a null set. For an integer period it becomes
\begin{equation}
 Q_{L,x}(f)=\frac1{pL}\left\{
 p^2A(f)+2\sum_{1\le j<L}
 \left|\sum_{a=1}^pe^{2\pi ijx_a/L}\right|^2C_{j/L}(f)\right\}.
 \label{eq:integer-period-energy}
\end{equation}
This uses Hilbert-space Parseval rather than a pointwise
Poisson-summation or Fourier-decay assertion.

Apply \eqref{eq:joint-affine-floor} to $M$ consecutive periods.
Their span divided by their number of points tends to $L/p$.
Their off-diagonal square energy per point tends to $Q_{L,x}(f)-1$:
each fixed displacement occurs with a weight tending to one and
bounded by one, while Parseval supplies the summable majorant.
Consequently \eqref{eq:defect-frobenius} gives
\begin{equation}
 a\le\sum_iw_i\{Q_{L,x}(f_i)-1\}+dL/p.
 \label{eq:periodic-floor-test}
\end{equation}
We may average finitely many such inequalities with nonnegative
weights. This does not require concatenating the periodic states
or realizing their mixture as a single zero process.

\subsection{The residual identity and its quantifiers}

Fix $0<u<1$, put $q=1-u$, and choose a finite mixture of periodic
states whose average gap is $1/q$. Let $Q_q(f)$ be its mixture
of full per-point energies. For a row define
\begin{equation}
 \mathcal F_q(f)=R(f)-qQ_q(f),\qquad
 \sigma=\sum_iw_i\{Q_q(f_i)-1\}+d/q-a.
 \label{eq:periodic-slack}
\end{equation}
Equation \eqref{eq:periodic-floor-test} gives $\sigma\ge0$.
Let $\operatorname{Res}$ be the left side of
\eqref{eq:joint-master-row} minus its right side. Substitution gives
the exact identity, valid also at $e=u$,
\begin{equation}
 \operatorname{Res}
 =\sum_iw_i\left\{\mathcal F_q(f_i)-\Gamma_i u
                    +\frac{\Gamma_i^2}{8}(u-e)\right\}+q\sigma.
 \label{eq:residual-uncompleted}
\end{equation}
For $0\le e<u$, completion of each scalar square gives
\begin{equation}
 \begin{split}
 \operatorname{Res}={}&
 \sum_iw_i\mathcal F_q(f_i)-\frac{2u^2}{u-e}\\
 &+\frac{u-e}{8}\sum_iw_i
       \left(\Gamma_i-\frac{4u}{u-e}\right)^2+q\sigma.
 \end{split}
 \label{eq:general-periodic-residual}
\end{equation}
The optimizing threshold $4u/(u-e)\ge4$ is permitted.
At $e=0$ its contribution is $(u/8)\sum_iw_i(\Gamma_i-4)^2$.

Suppose $\mathcal F_q$ has a minimizer $f_q^\dagger$
on the real affine space $\int f=1$, with minimum $m_q$, and
\[
 \mathcal F_q(h)\ge\kappa_q\|h\|_2^2
 \quad\left(\int h=0\right),\qquad\kappa_q>0.
\]
Stationarity makes the linear cross term vanish:
\begin{equation}
 \mathcal F_q(f)=m_q+\mathcal F_q(f-f_q^\dagger).
 \label{eq:general-variational-defect}
\end{equation}
Thus every row is satisfied at the same count point whenever
\begin{equation}
 m_q\ge\frac{2u^2}{u-e}.
 \label{eq:general-count-frontier}
\end{equation}
The residual retains exactly the profile departure, threshold
departure and periodic-test slack; its profile contribution is at
least $\kappa_q\sum_iw_i\|f_i-f_q^\dagger\|_2^2$.

The common point makes the obstruction simultaneous.
The mixture and $(u,e)$ are fixed before a row is chosen; only
the nonnegative slack $\sigma$ depends on that row.
Arbitrarily many already valid limiting rows can therefore be
retained together. This does not establish height-uniform control
for a growing family of tests or growing thresholds.

No division by $u-e$ is made at $e=u$:
\eqref{eq:residual-uncompleted} remains the appropriate formula.
A point with $e=u>0$ cannot survive the complete collection of
all-threshold rows. Indeed fix one admissible profile and use
$D_{\Gamma/2-1}\ge0$, so $a=d=B=0$. Its row requires
$R(f)\ge1-u+\Gamma u$ for every fixed $\Gamma\ge2$, which is
impossible as $\Gamma$ increases. The endpoint $u=e=0$
is tested directly in \eqref{eq:joint-master-row} and is not
represented by the quotient in \eqref{eq:general-count-frontier}.

\section{An explicit two-lattice variational frontier}
\label{sec:two-lattice}

The simplest density-balanced mixture gives an exact two-count
model and a quantitative rigidity identity. Its simple-critical
endpoint is less restrictive numerically than the construction
of Section~\ref{sec:twelve-fibre}, but the whole frontier has
a closed form and applies to the union count.

\subsection{The quadratic form and its coercivity}

The unit lattice and the lattice of spacing two have full energies
$A(f)$ and $A(f)/2+C_{1/2}(f)$, respectively.
For $1/8\le u\le1/3$, put $q=1-u$ and give the unit lattice
weight $\theta=2-1/q$. Then
$\theta+2(1-\theta)=1/q$, and the quadratic form is
\begin{equation}
 F_u(f)=J(f)+\frac{3u}{2}A(f)-uC_{1/2}(f).
 \label{eq:two-lattice-form}
\end{equation}
Its unique normalized real minimizer will be positive and even,
with minimum
\begin{equation}
 m(u)=\frac12+u\{\alpha\cot\alpha-2\beta\tan\beta\},
 \qquad
 \alpha=\frac1{2\sqrt{2u}},\quad
 \beta=\frac1{4\sqrt u}.
 \label{eq:two-lattice-minimum}
\end{equation}
The estimate needed down to $u=1/8$ is
\begin{equation}
 F_u(h)\ge\left(u-\frac18+\frac1{1000}\right)\|h\|_2^2
 \qquad\left(\int_Ih=0\right).
 \label{eq:two-lattice-coercivity}
\end{equation}

To prove it put $H(s)=\int_{-1/2}^sh(t)\dd t$.
Both endpoint values of $H$ vanish. For
$g_h(s)=\int|s-t|h(t)\dd t$, one has $g_h'=2H$;
integration by parts gives
\begin{equation}
 J(h)=-2\int_IH(s)^2\dd s.
 \label{eq:primitive-negative-energy}
\end{equation}
Let $B$ be half the measure-preserving half-interval swap:
$(Bh)(s)=h(s+1/2)/2$ for $s<0$ and
$(Bh)(s)=h(s-1/2)/2$ for $s>0$.
Then $B=B^*$, $B^2=I/4$, and
$C_{1/2}(h)=\langle h,Bh\rangle$.
The forms $A,J,C_{1/2}$ preserve reflection parity,
so their even--odd cross terms vanish.

If $h$ is even and zero-mean, its primitive also vanishes at zero.
Poincare's inequality on each half-interval gives
$J(h)\ge-\|h\|_2^2/(2\pi^2)$.
Since $C_{1/2}(h)\le A(h)/2$,
\[
 F_{1/8}(h)\ge
 \left(\frac18-\frac1{2\pi^2}\right)\|h\|_2^2
 >\frac1{1000}\|h\|_2^2.
\]
For odd $h$, write $h=t\psi+v$ with
$\psi(s)=\sqrt2\sin(\pi s)$ and $v\perp\psi$ odd.
The odd Neumann basis consists of
$\sqrt2\sin((2j+1)\pi s)$, $j\ge0$.
Its primitives are orthogonal Dirichlet modes of squared norm
$1/((2j+1)^2\pi^2)$. Hence
\[
 J(h)\ge-\frac2{\pi^2}t^2-\frac2{9\pi^2}\|v\|_2^2.
\]
Direct integration and $B^2=I/4$ give
\[
 \langle\psi,B\psi\rangle=-1/\pi,\qquad
 \|B\psi+\psi/\pi\|_2^2=1/4-1/\pi^2.
\]
Using $\langle v,Bv\rangle\le\|v\|_2^2/2$ now yields
\[
 F_{1/8}(h)\ge a_0t^2+b_0\|v\|_2^2-2c_0|t|\|v\|_2,
\]
where
\[
 a_0=\frac3{16}+\frac1{8\pi}-\frac2{\pi^2},\quad
 b_0=\frac18-\frac2{9\pi^2},\quad
 c_0^2=\frac1{64}\left(\frac14-\frac1{\pi^2}\right).
\]
For a rational positivity check put $\pi_L=157/50$, $\pi_U=22/7$,
\[
 a_L=\frac3{16}+\frac1{8\pi_U}-\frac2{\pi_L^2},\quad
 b_L=\frac18-\frac2{9\pi_L^2},\quad
 c_U^2=\frac1{64}\left(\frac14-\frac1{\pi_U^2}\right).
\]
Both $a_L-1/1000$ and $b_L-1/1000$ are positive, and
\[
 (a_L-1/1000)(b_L-1/1000)-c_U^2
 =\frac{17310471944257}{330823607944500000}>0.
\]
The two-by-two form is therefore at least
$(t^2+\|v\|_2^2)/1000$. This proves the estimate at $u=1/8$.
Finally
\[
 F_u(h)-F_{1/8}(h)
 =(u-1/8)\{3A(h)/2-C_{1/2}(h)\}
 \ge(u-1/8)A(h),
\]
which proves \eqref{eq:two-lattice-coercivity}.

\subsection{The minimizer and its exact defect}

With $\alpha,\beta$ as in \eqref{eq:two-lattice-minimum}, set
$p=\alpha/\sin\alpha$ and $z=-\beta/\cos\beta$, and define
\begin{equation}
 f_u(s)=p\cos\bigl(\sqrt{2/u}(|s|-1/4)\bigr)
       +z\sin\bigl((|s|-1/4)/\sqrt u\bigr).
 \label{eq:two-lattice-profile}
\end{equation}
The sine term integrates to zero on each half-interval, and
$4p\sin\alpha/\sqrt{2/u}=1$, so $\int f_u=1$.
For $g(s)=\int|s-t|f_u(t)\dd t$ the Euler expression on
$0<s<1/2$ is
\[
 L(s)=3u f_u(s)-u f_u(1/2-s)+2g(s).
\]
The cosine component is symmetric about $1/4$ and solves
$2u w''+4w=0$; the sine component is antisymmetric and solves
$4u w''+4w=0$. As $g''=2f_u$, these imply $L''=0$.
The endpoint derivatives are
$f_u'(0+)=1/(4u)$ and $f_u'(1/2-)=-3/(4u)$.
Together with $g'(0)=0$ they give $L'(0+)=0$.
Thus $L$ is constant on the positive half. Reflection makes the
same constant valid on the negative half, establishing
stationarity against all real zero-mean perturbations.

Coercivity gives unique global minimality. At $s=1/2$,
evenness and normalization imply $g(1/2)=1/2$.
Multiplication of the Euler equation by $f_u$ and integration
show that its constant is $2F_u(f_u)$.
Endpoint evaluation gives \eqref{eq:two-lattice-minimum}.
Stationarity also proves, for every real normalized $f$,
\begin{equation}
 F_u(f)=m(u)+F_u(f-f_u).
 \label{eq:two-lattice-exact-defect}
\end{equation}

To prove positivity, write $x=s-1/4$ on the positive half.
For $x\ge0$ both derivative terms in
\eqref{eq:two-lattice-profile} are nonpositive, and
$f_u(1/4-x)\ge f_u(1/4+x)$. The minimum is consequently
\[
 f_u(1/2)=\alpha\cot\alpha-\beta\tan\beta.
\]
At $u=1/8$ the arguments are $1$ and $1/\sqrt2$.
Alternating Taylor bounds give
\[
 \cos1\ge389/720,\quad \sin1\le101/120,\quad
 \frac{\sin(1/\sqrt2)}{1/\sqrt2}\le147/160,\quad
 \cos(1/\sqrt2)\ge4379/5760.
\]
Therefore
\[
 f_{1/8}(1/2)\ge
 \frac{389}{606}-\frac{2646}{4379}
 =\frac{99955}{2653674}>0.
\]
On $(0,\pi/2)$, $x\cot x$ decreases and $x\tan x$ increases.
Increasing $u$ decreases both arguments, so the endpoint
remains positive. The signed affine minimum is thus also the
minimum over nonnegative profiles.

\subsection{The frontier, its rigidity and the union witness}

The two lattice tests give nonnegative slacks
\[
 \sigma_1=\overline A-1+d-a,\qquad
 \sigma_2=\overline A/2+\overline C_{1/2}-1+2d-a.
\]
Since $q\theta=1-2u$ and $q(1-\theta)=u$,
$q\sigma=(1-2u)\sigma_1+u\sigma_2$.
The general residual is therefore
\begin{equation}
 \begin{split}
 \operatorname{Res}={}&m(u)-\frac{2u^2}{u-e}
       +\sum_iw_iF_u(f_i-f_u)\\
 &+\frac{u-e}{8}\sum_iw_i
       \left(\Gamma_i-\frac{4u}{u-e}\right)^2
       +(1-2u)\sigma_1+u\sigma_2.
 \end{split}
 \label{eq:two-lattice-frontier-residual}
\end{equation}
All terms after the first are nonnegative.
The positive profile gives $m(u)>0$, since $J(f_u)>0$ and
$3A(f_u)/2-C_{1/2}(f_u)\ge A(f_u)>0$.
The exact sufficient frontier is
\begin{equation}
 1/8\le u\le1/3,\quad0\le e<u,\qquad
 e\le E(u):=u-\frac{2u^2}{m(u)}.
 \label{eq:two-count-frontier}
\end{equation}
The profile term costs at least
$(u-1/8+1/1000)\sum_iw_i\|f_i-f_u\|_2^2$.
Near equality at the frontier controls the weighted squared departures
from the same profile and threshold, together with the lattice-test
slacks. Individual control requires a lower bound for that profile's weight.

Exactness refers to the enlarged dual problem constrained only
by these two lattice tests. To verify it, use one profile $f_u$,
write $E_1=A(f_u)-1$ and
$E_2=A(f_u)/2+C_{1/2}(f_u)-1$, and choose
\[
 d=E_1-E_2\ge0,\qquad a=2E_1-E_2,\qquad
 \Gamma=4u/(u-e).
\]
Both lattice tests are equalities and every residual term after
the first vanishes. No actual universal Gram floor with those
coefficients has been asserted; other periodic states and
clipping can impose further restrictions.

The function $m$ is smooth by its formula and concave as an
infimum of affine functions of $u$. Stationarity gives
$m'(u)=3A(f_u)/2-C_{1/2}(f_u)$, so $m-u m'=J(f_u)>0$.
Differentiation yields
\[
 \left(\frac{u^2}{m}\right)''
 =\frac{2(m-u m')^2}{m^3}-\frac{u^2m''}{m^2}>0.
\]
Thus $E$ is strictly concave on the stated interval.
The following consequence needs no exact maximizer:
\begin{equation}
 u_{\rm U}=\frac{13711914}{10^8},\qquad
 e_{\rm U}=\frac{5251083389778}{10^{14}},\qquad
 m(u_{\rm U})>\frac{2u_{\rm U}^2}{u_{\rm U}-e_{\rm U}}.
 \label{eq:union-common-point}
\end{equation}

For a finite rational verification, let $z=1/(16u)$ and put
\[
 C_n(z)=\sum_{k=0}^n\frac{(-1)^kz^k}{(2k)!},\qquad
 S_n(z)=\sum_{k=0}^n\frac{(-1)^kz^k}{(2k+1)!}.
\]
For $0<2z\le1$, the index-15 sums bound the corresponding
cosine and normalized sine below, and the index-16 sums bound
them above; all denominator lower bounds are positive.
It follows that
\[
 m(u)\ge m_L(u):=\frac12+
 u\left\{\frac{C_{15}(2z)}{S_{16}(2z)}
             -2z\frac{S_{16}(z)}{C_{15}(z)}\right\}.
\]
Substitution of the rationals in \eqref{eq:union-common-point}
and positive-denominator cross multiplication give
\[
 m_L(u_{\rm U})
 -\frac{2u_{\rm U}^2}{u_{\rm U}-e_{\rm U}}>10^{-14}.
\]
This specifies the certificate by finite rational expressions;
no quadrature or optimization is a premise.

The union complement consists of multiple nonreal pairs.
A pair of multiplicity $j\ge2$ has excluded mass $2j$
and excess $2(j-1)$, so $N-N_{\rm s\cup0}\le2eN$.
At \eqref{eq:union-common-point}, equality is compatible with
the count variables: assign mass $1-u_{\rm U}$ to simple real
locations, $u_{\rm U}-2e_{\rm U}$ to simple nonreal pairs, and
$2e_{\rm U}$ to double nonreal pairs. These masses are
nonnegative. A double pair has mass four and excess two;
multiplying by a common integer denominator gives integral
populations with precisely the specified excess. The retained
union relaxation therefore cannot force a lower proportion above
\begin{equation}
 1-2e_{\rm U}
 =\frac{22374458305111}{25000000000000}
 =0.89497833220444.
 \label{eq:union-method-ceiling}
\end{equation}
This is a feasible count population, not a construction of
exponential vectors realizing every analytic energy.

The zero-excess endpoint is the root $r_*\in[8/25,1/3]$
of $m(r_*)=2r_*$. On that smaller interval
$\alpha^2\le25/64$ and $\beta^2\le25/128$; elementary sine and
cosine bounds imply
\[
 \|f_u\|_\infty\le\frac{384}{359}+\frac{50}{231}<\frac43.
\]
As $f_u\ge0$ and $\int f_u=1$, this gives $m'(u)<2$.
Thus $m(u)-2u$ decreases strictly there.
The same alternating expressions, with indices 9 and 10, give
opposite signs at $32349386085/10^{11}$ and
$32349386086/10^{11}$. Hence
\[
 0.67650613914<1-r_*<0.67650613915.
\]
This is the simple-count endpoint of the two-lattice frontier;
the next construction gives a stronger numerical obstruction.

\section{A twelve-fibre obstruction for all profiles}
\label{sec:twelve-fibre}

Gebendorfer \cite[Theorem 1.2 and Sections 5--6]{GebendorferCap}
used two periodic configurations to obtain the ceiling $0.674034$
for a normalized nonnegative 25-mode profile family, with 24 free
coefficients. The periodic configurations and their density
balancing are his inputs. We solve the associated variational
problem on the full affine $L^2$ space at the slightly larger value
$0.67404$. We also give a positive profile showing why the
original finite-mode test cannot be extended unchanged.

\subsection{The quadratic problem and its coercivity}

The two configurations are
\[
 A_{\rm per}:\ (0,\alpha_0)\pmod3,\quad \alpha_0=5189/5000,
 \qquad
 B_{\rm per}:\ (0,\gamma_0,2\gamma_0)\pmod4,\quad
 \gamma_0=10289/10000.
\]
By \eqref{eq:integer-period-energy}, their full per-point energies are
\begin{equation}
 \begin{split}
 Q_A(f)&=\frac16\left[4A(f)+2\sum_{j=1}^2
 \{2+2\cos(2\pi j\alpha_0/3)\}C_{j/3}(f)\right],\\
 Q_B(f)&=\frac1{12}\left[9A(f)+2\sum_{j=1}^3
 \{3+4\cos(\pi j\gamma_0/2)+2\cos(\pi j\gamma_0)\}
 C_{j/4}(f)\right].
 \end{split}
 \label{eq:twelve-periodic-energies}
\end{equation}
For $2/3<q<3/4$, give the second configuration weight
$\lambda_q=9-6/q$. Then
$(1-\lambda_q)3/2+\lambda_q4/3=1/q$.
Writing $Q_q=(1-\lambda_q)Q_A+\lambda_qQ_B$, we have
\begin{equation}
 \mathcal F_q(f)=R(f)-qQ_q(f)
 =R(f)-(6-8q)Q_A(f)-(9q-6)Q_B(f).
 \label{eq:twelve-form}
\end{equation}

Every correlation shift is a multiple of $1/12$.
Put $\ell=1/12$, $x_j=-1/2+j\ell$, and
$F(x)=(f(x_0+x),\ldots,f(x_{11}+x))^t$ for $0<x<\ell$.
Let $T_k$ have ones at $(j,j+k)$, $0\le j<12-k$, and zeros
elsewhere. These are truncated shifts, with no periodic
identification of cell endpoints. If
$Q_q=c_0A+\sum_{k\in\{3,4,6,8,9\}}c_kC_{k/12}$, then
\begin{equation}
 \begin{split}
 c_0&=(1-\lambda_q)\frac23+\lambda_q\frac34,\\
 c_{4j}&=\frac{1-\lambda_q}{3}
       \{2+2\cos(2\pi j\alpha_0/3)\},\qquad j=1,2,\\
 c_{3j}&=\frac{\lambda_q}{6}
       \{3+4\cos(\pi j\gamma_0/2)+2\cos(\pi j\gamma_0)\},
       \qquad j=1,2,3.
 \end{split}
 \label{eq:twelve-coefficients}
\end{equation}
The two index sets in the last lines are disjoint.
Since $C_{k/12}(f)=\int_0^\ell F(x)^tT_kF(x)\dd x$,
\begin{equation}
 \begin{split}
 M_q&=(1-qc_0)I-\frac q2\sum_{k\in\{3,4,6,8,9\}}
                         c_k(T_k+T_k^t),\\
 \mathcal F_q(f)&=\int_0^\ell F(x)^tM_qF(x)\dd x+J(f).
 \end{split}
 \label{eq:twelve-matrix}
\end{equation}

Set
\begin{equation}
 q_{25}=\frac{337017}{500000},\qquad
 q_{\rm c}=\frac{16851}{25000}.
 \label{eq:twelve-test-values}
\end{equation}
The explicit rational enclosures in Appendix~\ref{app:twelve-fibre}
give, at both values, every diagonal-dominance margin of $M_q$
greater than $77/250$. Symmetry therefore gives
$M_q>(77/250)I$. For a zero-mean perturbation,
\eqref{eq:primitive-negative-energy} and the Dirichlet
Poincare inequality on the whole interval give
$J(h)\ge-2\|h\|_2^2/\pi^2$. Consequently
\begin{equation}
 \mathcal F_q(h)>
 \left(\frac{77}{250}-\frac2{(157/50)^2}\right)\|h\|_2^2
 >\frac1{10}\|h\|_2^2\qquad(h\ne0,\ \int h=0).
 \label{eq:twelve-coercivity}
\end{equation}
All forms are bounded on $L^2(I)$. Translate the affine space
by any normalized function and restrict to the closed
zero-mean subspace. The bounded positive operator representing
the coercive bilinear form is invertible there, and its inverse
solves the stationarity equation. Thus a unique affine
minimizer exists. This controls the full space, with no
finite-mode remainder omitted.

\subsection{A posteriori certification on the affine space}

For any normalized real candidate $f_0$, put
$g_0(s)=\int_I|s-t|f_0(t)\dd t$ and let
\[
 r=M_qf_0+g_0-c\one
\]
for any real $c$, interpreting the multiplication in fibres.
If $m_q=\min_{\int f=1}\mathcal F_q(f)$, then
\begin{equation}
 \mathcal F_q(f_0)-10\|r\|_2^2
 \le m_q\le\mathcal F_q(f_0).
 \label{eq:twelve-residual-certificate}
\end{equation}
Indeed, for $\int h=0$,
\[
 \mathcal F_q(f_0+h)
 =\mathcal F_q(f_0)+2\langle r,h\rangle+\mathcal F_q(h)
 \ge\mathcal F_q(f_0)+2\langle r,h\rangle+\tfrac1{10}\|h\|_2^2.
\]
Completion of this Hilbert-space square proves the lower bound;
the candidate itself proves the upper bound. The constant $c$
disappears from the cross term because $\int h=0$.

Appendix~\ref{app:twelve-fibre} defines one even, normalized,
piecewise polynomial $f_0>7/10$ using 42 integer coefficients.
Its evaluation is rational after seven cosine values have been
enclosed by the displayed intervals. The resulting coarse
enclosures, sufficient for the proof, are
\begin{equation}
 \begin{split}
 \frac{651926252579}{10^{12}}
 &<\mathcal F_{q_{25}}(f_0)
   <\frac{651926252580}{10^{12}},\\
 \frac{651920180779}{10^{12}}
 &<\mathcal F_{q_{\rm c}}(f_0)
   <\frac{651920180780}{10^{12}},\\
 \left\|M_{q_{\rm c}}f_0+g_0-
       \frac{651920180779}{10^{12}}\one\right\|_\infty
 &<\frac4{10^{12}}.
 \end{split}
 \label{eq:twelve-certificate-values}
\end{equation}
Since $I$ has length one, the squared $L^2$ residual is less
than $16/10^{24}$. It follows that
\begin{equation}
 m_{q_{\rm c}}>
 \frac{651920180779}{10^{12}}-\frac{160}{10^{24}}
 >2-2q_{\rm c}+\frac{18}{10^8},
 \qquad 2-2q_{\rm c}=\frac{8149}{12500}.
 \label{eq:twelve-positive-margin}
\end{equation}
On the other hand the first line of
\eqref{eq:twelve-certificate-values} gives
\begin{equation}
 \mathcal F_{q_{25}}(f_0)
 <2-2q_{25}-\frac1{200000}.
 \label{eq:old-periodic-test-fails}
\end{equation}
Thus this periodic test does not exclude its original value over
the larger profile space. This neither contradicts the source's
finite-dimensional theorem nor excludes a different all-profile
test at the 25-mode value.

\subsection{Positivity of the true minimizer}

The signed minimization already suffices for the obstruction.
It is also the exact minimization over nonnegative profiles.
Let $f^\dagger$ be the minimizer at $q_{\rm c}$ and put
$c_0=\mathcal F_{q_{\rm c}}(f_0)$. The displayed enclosures give
\[
 \|M_{q_{\rm c}}f_0+g_0-c_0\one\|_\infty
 <5\cdot10^{-12}<\varepsilon,\qquad\varepsilon=10^{-11}.
\]
Call this centered residual $r_0$ and set $h=f^\dagger-f_0$.
The minimizing equation is
$M_{q_{\rm c}}f^\dagger+g_{f^\dagger}=m_{q_{\rm c}}\one$:
multiplication by $f^\dagger$ and integration identify its
multiplier with the minimum. Subtract the candidate equation
and pair with $h$ to obtain
$\mathcal F_{q_{\rm c}}(h)=-\langle r_0,h\rangle$.
Coercivity and \eqref{eq:twelve-residual-certificate} imply
\[
 \|h\|_2\le10\varepsilon,\qquad
 0\le c_0-m_{q_{\rm c}}\le10\varepsilon^2.
\]
Because $|s-t|\le1$, $\|g_h\|_\infty\le\|h\|_2$.
Strict row diagonal dominance gives
$\|M_{q_{\rm c}}^{-1}\|_{\ell^\infty\to\ell^\infty}
 <250/77<4$: apply a row at a coordinate of maximum modulus
in a finite-vector equation. The difference of the Euler
equations therefore gives
\[
 \|h\|_\infty<
 4(\varepsilon+10\varepsilon+10\varepsilon^2)<10^{-9}.
\]
The equation supplies this bounded representative fibre by fibre;
boundedness was not presumed. Since $f_0>7/10$,
$f^\dagger$ is positive almost everywhere.
Reflection invariance and uniqueness also make it even.

For completeness there is an exact differential characterization,
although certification did not require its numerical inversion.
Put $G_j(x)=g_{f^\dagger}(x_j+x)$ and
$V=G-m_{q_{\rm c}}\one$. Then
\[
 V''+2M_{q_{\rm c}}^{-1}V=0,\qquad
 F=-M_{q_{\rm c}}^{-1}V.
\]
The 25 scalar conditions are continuity of $G,G'$ at the eleven
internal junctions, $G'_0(0)=-1$, $G'_{11}(\ell)=1$, and
$G_0(0)+G_{11}(\ell)=1$.
The endpoint derivatives enforce normalization through $g''=2f$.
The endpoint sum follows from
$|t+1/2|+|1/2-t|=1$.
Conversely these conditions identify $G$ with the integral
potential: their difference is globally affine, with zero
endpoint derivatives and zero endpoint sum.
Continuity is imposed on the potential and its derivative,
not on the profile; the latter may jump at cell boundaries.

All forms in \eqref{eq:structural-forms} are continuous in $L^2$:
shifts are contractions, and $|s-t|$ is a bounded integral
kernel on an interval of finite measure.
The positive even $f_0$ and $f^\dagger$ can be approximated by
profiles whose square roots are smooth and supported strictly
inside $I$, as required in Section~\ref{sec:interface}.
For either $f$, first approximate $\sqrt f$ in $L^4$ by
nonnegative even $\eta_j\in C_c^\infty((-1/2,1/2))$, using
symmetric cutoffs and even nonnegative mollifiers. The estimate
\[
 \|\eta_j^2-f\|_2
 \le\|\eta_j-\sqrt f\|_4
       \bigl(\|\eta_j\|_4+\|\sqrt f\|_4\bigr)
 \longrightarrow0
\]
shows that $\eta_j^2/\int\eta_j^2\to f$ in $L^2$.
These normalized profiles have the required smooth square roots.
The strict inequality \eqref{eq:old-periodic-test-fails} persists
for a sufficiently close approximation. The variational infimum
at $q_{\rm c}$ is unchanged on this dense admissible class,
although its minimizer need not belong to the smooth class.

\subsection{The simultaneous method ceiling}

Apply \eqref{eq:general-periodic-residual} at the single count point
$u=1-q_{\rm c}$, $e=0$. Stationarity gives for every covered row
\begin{equation}
 \begin{split}
 \operatorname{Res}={}&
 \delta+\sum_iw_i\mathcal F_{q_{\rm c}}(f_i-f^\dagger)\\
 &+\frac{1-q_{\rm c}}8\sum_iw_i(\Gamma_i-4)^2
       +q_{\rm c}\sigma,\qquad
 \delta=m_{q_{\rm c}}-(2-2q_{\rm c})>\frac{18}{10^8}.
 \end{split}
 \label{eq:twelve-joint-residual}
\end{equation}
The profile term is at least
$\frac1{10}\sum_iw_i\|f_i-f^\dagger\|_2^2$.
The same point therefore satisfies every joint-affine row,
including arbitrary individual clipping thresholds.
Its simple-real proportion is
\begin{equation}
 q_{\rm c}=0.67404,\qquad
 \text{giving the method ceiling }67.404\%.
 \label{eq:all-profile-method-ceiling}
\end{equation}
All omitted mass can be assigned to simple nonreal pairs.
Every location is then simple and the distinct, simple-anywhere
and union counts all equal $N$. Lower bounds for those statistics
do not remove this count witness.

The comparison with earlier ceilings has two dimensions.
Gebendorfer's smaller $0.674034$ value applies on his
finite-dimensional family; the present result covers every
normalized $L^2$ profile in the specified joint-affine interface.
Knausg{\aa}rd's fixed second-window bounds
\cite[Proposition 7.3]{Knausgard2026} are smaller still:
$0.67357$ for seven index lags and $0.67383$ for all pairs.
His periodic-law result
\cite[Theorem 7.7 and Corollary 7.9]{Knausgard2026}, following
\cite[Section 7.2]{AF}, treats a broader trace--energy construction
with explicit regularity and variation costs.
The present proof neither replaces every such construction nor
identifies the optimal periodic mixture. It establishes the
all-profile obstruction and quantitative residual for the rows
defined in Section~\ref{sec:structural}.


\section{Interaction retained and information still available}

The constructions above separate three uses of a quadratic zero energy.
First, separation can make clipping exact on a large block. For the
explicit cosine profile this retains the full local reward and yields
\eqref{eq:intro-main}. Second, local slack can control geometry rather
than merely improve a scalar constant: the phase comparison bounds
the departure of adjacent gaps from integral spacing, with an explicit
price for arbitrarily long physical spans. Third, periodic configurations
test the information content of affine deductions simultaneously over
profiles and thresholds.

The variational results make the third distinction quantitative.
Coercivity identifies the unique optimal profile for a prescribed
periodic mixture and measures the cost of departing from it. The
common residual identifies the independent cost of a nonoptimal
threshold and the slack in the periodic tests. These are statements
about the data retained by the deduction. They explain why further
information about the same operator can matter even when an affine
count has already been optimized.

The numerical proofs use two kinds of finite input. The eight-point
floor for the main profile has full analytic enclosure and finite
certificate data in the accompanying formal source. The labelled,
third-profile and compensated-frame inequalities used elsewhere are
the computer-assisted results cited at their statements; their
published certificate computations are not replaced by the small
arithmetic checks accompanying this paper. Those checks reproduce
the displayed energy enclosures, table values, phase-polynomial
bounds and variational residual certificates. In particular, the
additional variational and multiplicity deductions are not asserted
to be covered by the final exported statement of the supplied Lean
file.

\section*{acknowledgments}
The author declares the substantial use of consumer-grade Gemini-3 and GPT-6 models for literature search, exploring and developing ideas, codes and proofs, drafting and LaTeX typesetting; the author takes full and sole responsibility for all contents of this work.


\section*{Data availability}
Numerical results reported here are reproducible from the verification
scripts and certificate data in the repository. Lean code and verification details are also included. - Repository - 
\href{https://github.com/arun-chandru/riemann-zeros}{github.com/arun-chandru/riemann-zeros}.



\appendix
\section{The imported eight-point data and their algebraic consequences}
\label{app:local-data}

This appendix specifies the local input for the second profile
\eqref{eq:profile-two}. The whole-domain inequality
\eqref{eq:local-floor} is the computer-assisted input of
\cite[Proposition 5.3 and Tables 10--11]{Knausgard2026}, at
$\theta=17/20$ and $\alpha=627/10^6$, with reward $a=8722/10^6$.
The tables below fix that premise completely. The finite capacity,
potential, and slack calculations following them are separate elementary
consequences. Listing the data is not a replay of the source's continuum
search or its computation-dependent formalisation.

All entries in the next two tables are integer numerators with denominator
$10^6$. In the first table the row is $r$, the column is $s$, and the entry
is $10^6W_{rs}$; a dash is outside the range $0\le r<s\le7$.
\begin{center}
\small\setlength{\tabcolsep}{4pt}
\begin{tabular}{r*{7}{r}}
\toprule
$r\backslash s$&1&2&3&4&5&6&7\\\midrule
0&159939&0&280438&143766&709387&1000000&2000000\\
1&---&264496&385569&340591&856233&581225&1000000\\
2&---&---&389188&614430&757940&856233&709387\\
3&---&---&---&372750&614430&340591&143766\\
4&---&---&---&---&389188&385569&280438\\
5&---&---&---&---&---&264496&0\\
6&---&---&---&---&---&---&159939\\\bottomrule
\end{tabular}
\end{center}
The corrections are continuous piecewise-linear functions of a gap.
Their values at the fourteen knots are as follows; all seven functions
vanish outside $[3/4,19/8]$.
\begin{center}
\small\setlength{\tabcolsep}{5pt}
\begin{tabular}{r*{7}{r}}
\toprule
$t$&$10^6B_0(t)$&$10^6B_1(t)$&$10^6B_2(t)$&$10^6B_3(t)$&
$10^6B_4(t)$&$10^6B_5(t)$&$10^6B_6(t)$\\\midrule
$3/4$&0&0&0&0&0&0&0\\
$7/8$&647&0&-425&-444&-425&0&647\\
$1$&555&91&-425&-442&-425&91&555\\
$9/8$&518&45&-342&-442&-342&45&518\\
$5/4$&647&0&-425&-444&-425&0&647\\
$11/8$&647&0&-425&-444&-425&0&647\\
$3/2$&647&0&-425&-444&-425&0&647\\
$13/8$&0&0&221&-442&221&0&0\\
$7/4$&0&0&221&-442&221&0&0\\
$15/8$&270&86&-275&-162&-275&86&270\\
$2$&325&58&-244&-278&-244&58&325\\
$17/8$&0&0&-425&850&-425&0&0\\
$9/4$&647&0&-425&-444&-425&0&647\\
$19/8$&0&0&0&0&0&0&0\\\bottomrule
\end{tabular}
\end{center}
Thus $(B_0,\ldots,B_6)=(u_0,u_1,u_2,-2(u_0+u_1+u_2),
u_2,u_1,u_0)$, where $u_j=B_j$ for $j=0,1,2$. In particular
$\sum_{r=0}^6B_r(t)=0$ identically, including between knots.

The seven lag-capacity sums, in the same integer units, are
\begin{equation}
 10^6\sum_{r=0}^{7-h}W_{r,r+h}
 =(1999996,1999998,1999998,1999998,1999999,2000000,2000000)_h,
 \quad 1\le h\le7.
 \label{eq:appendix-capacities}
\end{equation}
Every sum is at most $2\cdot10^6$. A pair in a translated configuration
can occur at most once at each admissible relative position, so these
are exactly the capacity bounds needed to charge its total weight to
the ordered energy with coefficient two.

For completeness, the bounded potential can also be read directly from
the table. Set $C_r=\sum_{s=0}^rB_s$ for $0\le r\le5$ and
\[
 \Phi(g_0,\ldots,g_5)=\sum_{r=0}^5C_r(g_r).
\]
The ranges of $10^6C_r$ are respectively
\[
 [0,647],\quad[0,647],\quad[-425,222],\quad[-222,425],
 \quad[-647,0],\quad[-647,0].
\]
These ranges include the zero values outside the knot interval.
Linear interpolation cannot create a new extremum. Each range has
length $647$, and therefore
\begin{equation}
 \osc\Phi\le\Omega:=\frac{3882}{10^6},\qquad
 \sum_{r=0}^6B_r(g_r)
 =\Phi(g_0,\ldots,g_5)-\Phi(g_1,\ldots,g_6).
 \label{eq:appendix-potential}
\end{equation}
On summing complete frames the potential telescopes. The spread sums
to at most seven times the total spread, and the reward occurs $m-7$
times. Keeping the nonnegative frame slacks gives exactly
\eqref{eq:retained-slack}, with $B_0=7a+\Omega$; this scalar $B_0$
is distinct from the correction function $B_0(t)$ in the table.
For $m<8$ the same inequality holds with zero slack and empty sums.

The minima of the seven correction functions sum to
$-(425+444+425)/10^6=-1294/10^6$. Since all pair weights are
nonnegative, each frame slack is consequently at least
\[
 \alpha(y_{i+7}-y_i)-a-1294/10^6
 =\alpha(y_{i+7}-y_i)-10016/10^6.
\]
This calculation supplies the span price used in the new argument and
requires no positive margin for the source's accepted search boxes.

The associated separated-norm and close-pair inputs are also those of
\cite[Section 5]{Knausgard2026}. Its majorant is
\[
 g(u)=\frac{13}{20}
 \left\{\sinc\bigl(\theta(u-1/2)\bigr)
       +\sinc\bigl(\theta(u+1/2)\bigr)\right\}^2,
 \qquad\theta=17/20.
\]
The cited certificate proves $g\ge f$ on $I$ and
$K(t)^2\ge13707941/2^{28}$ on $|t|\le\theta$.
The Fourier transform of $g$ vanishes for $|\xi|\ge\theta$, and
\[
 \Lambda=\int_{\R}g
 =\frac{26}{17}(1+\sinc(17/20))<\frac{520}{289}<\frac95<2.
\]
Here $\sin(17\pi/20)=\sin(3\pi/20)<3\pi/20$.
The support and integral follow directly from the sinc representation;
the two continuum inequalities remain imported premises. These are the
premises used in the smoothing and matching argument, rather than an
inference about operator norms from entrywise approximation.

\section{Exact arithmetic for the phase comparison and profile energies}
\label{app:arithmetic}

The new finite calculations use rational arithmetic. This appendix gives
the recurrences, endpoints, and inequalities needed to reproduce them
without a floating-point grid or additional numerical data. The scalar
checks verify consequences of the stated imported premises; they do not
replace those premises.

\subsection{A common enclosure of \texorpdfstring{$\pi$}{pi} and the profile energies}

For $0<x<1$ put
\[
 A_n(x)=\sum_{k=0}^{n-1}\frac{(-1)^kx^{2k+1}}{2k+1},\qquad
 r_n(x)=\frac{x^{2n+1}}{2n+1}.
\]
For even $n$, the alternating-series remainder gives
$A_n(x)<\arctan x<A_n(x)+r_n(x)$. Machin's identity
$\pi=16\arctan(1/5)-4\arctan(1/239)$, with $n=24$ and $n=8$
respectively, proves by rational comparison
\begin{equation}
 \pi_-:=\frac{3141592653589793}{10^{15}}
 <\pi<
 \pi_+:=\frac{3141592653589794}{10^{15}}.
 \label{eq:appendix-pi}
\end{equation}
Explicitly the inner lower endpoint is
$16A_{24}(1/5)-4\{A_8(1/239)+r_8(1/239)\}$ and the inner upper
endpoint is $16\{A_{24}(1/5)+r_{24}(1/5)\}-4A_8(1/239)$;
they lie strictly between the displayed rational endpoints.

For the three integer-cosine profiles used above, write the energy as
$A+B/\pi^2$ using \eqref{eq:plus-cosine-energy}. The exact pairs are
\begin{center}
\small
\begin{tabular}{crr}
\toprule Profile&$A$&$B$\\\midrule
1&$1624643573617/1200000000000$&$-250664498438947/1200000000000000$\\
2&$4026042764683/3000000000000$&$-440830631846157799/3175200000000000000$\\
3&$4060769086579/3000000000000$&$-1000828214586143/4800000000000000$\\\bottomrule
\end{tabular}
\end{center}
All three values of $B$ are negative. Thus the energy lies between
$A+B/\pi_-^2$ and $A+B/\pi_+^2$, in that order. For convenient
outward endpoints, multiplication by $10^{15}$ gives
\begin{equation}
\begin{aligned}
 1332704958281710&<10^{15}R_1<1332704958281712,\\
 1327947272844555&<10^{15}R_2<1327947272844557,\\
 1332463633167058&<10^{15}R_3<1332463633167061.
\end{aligned}
 \label{eq:appendix-energy-enclosures}
\end{equation}
The first and third upper endpoints are below the caps in
\eqref{eq:plus-other-energy-caps}. The middle profile is used in the
phase comparison. All three enclosures follow from the same short
Machin calculation.

\subsection{The phase polynomial and its six rational intervals}

Use exactly the lags, radii, and weights in \eqref{eq:lag-data}.
Generate the Chebyshev coefficients by
$T_0=1$, $T_1=t$, $T_{l+1}=2tT_l-T_{l-1}$. Form the rational
numerator $N(t)=\sum_lw_l(1-T_l(t))-(1-t)$ in
\eqref{eq:phase-polynomial}. If $N(t)=\sum_{j=0}^{49}n_jt^j$,
then $\sum_jn_j=0$ and division by $1-t$ gives
\[
 P(t)=\sum_{k=0}^{48}p_kt^k,\qquad
 p_k=\sum_{j=0}^kn_j.
\]
The checks $n_{49}=-p_{48}$ and
$P(1)=\sum_lw_ll^2-1$ provide endpoint identities independent
of a sign test.

On any rational interval $[u,v]$, substitute $t=u+(v-u)x$.
Its power coefficients and degree-48 Bernstein coefficients are
\begin{equation}
 d_k=(v-u)^k\sum_{j=k}^{48}p_j\binom jk u^{j-k},\qquad
 b_s=\sum_{k=0}^s d_k\frac{\binom sk}{\binom{48}k},
 \quad0\le k,s\le48.
 \label{eq:appendix-bernstein}
\end{equation}
The Bernstein basis is
$\binom{48}s x^s(1-x)^{48-s}$; its entries are nonnegative and
sum to one on $[0,1]$. Direct use of
\eqref{eq:appendix-bernstein} gives, for all $49$ coefficients on
each interval, the strict lower bounds in Table~\ref{tab:phase-bounds}.
Every bound exceeds $1/10000$. These $294$ rational comparisons
prove the phase inequality on the entire circle. Equivalently, start
with the Bernstein coefficients on $[-1,1]$ and apply de Casteljau
midpoint subdivision to the six stated pieces; its updates are simply
$(b_j+b_{j+1})/2$. The independently implemented check instead
constructs $T_l$ from the closed formula
\[
 T_l(t)=\frac l2\sum_{j=0}^{\lfloor l/2\rfloor}
 \frac{(-1)^j(l-j-1)!}{j!(l-2j)!}(2t)^{l-2j}
 \quad(l\ge1)
\]
and uses direct interval conversion, avoiding the first implementation's
recurrence and subdivision.

The other rigidity checks are short finite sums. With $c_j=(-1)^ja_j$
and $z_0=2601/25$, the tail lower bound is
\[
 1+\sum_{c_j<0}\frac{c_jz_0}{z_0-j^2}+\sum_{c_j>0}c_j
 =\frac{3123731672835001}{3750959784000000}>\frac{104}{125}.
\]
For each lag the excess in the radius condition is exactly
\[
 R_l-\left\{q_ld_0/\alpha-(7q_l-l)\theta\right\}
 =\frac{16q_l}{627}>0.
\]
Using $\pi<22/7$, the remaining rational comparisons are
\begin{equation}
 \frac{196(22/7)^2}{8(104/125)^2}
 =\frac{1890625}{5408}<350,\qquad
 \sum_{l=12}^{49}\frac{196q_l}{2\alpha R_l^3}<112,
 \qquad
 \sum_{l=12}^{49}\frac{196l}{8R_l^2}<6.
 \label{eq:appendix-resource-checks}
\end{equation}
All denominators are positive. Together with the combinatorial proof in
Section~\ref{sec:rigidity}, these comparisons establish the displayed
energy, slack, and boundary prices.

\subsection{Primitive pressure values and scalar guards}

Integer derivatives can be reconstructed without trigonometric interval
arithmetic. Put $s(0)=0$ and $s(n)=(-1)^n/n$ for nonzero integers.
The sinc expansion gives
\[
 K'(k)=s(k)+\frac12\sum_{j=1}^{10}a_j\{s(k-j)+s(k+j)\},
 \qquad K(k)=a_k/2\quad(k=1,2).
\]
This yields the signed values in Section~\ref{sec:bootstrap}.
The profile absolute-coefficient sum is $210672/10^6<1$, and
the second-derivative bound has the explicit margin
\[
 4-\frac{(22/7)^2}{3}\left(1+\frac{210672}{10^6}\right)
 =\frac{31793}{2296875}>0.
\]

For $A_k=|K(k)|$, $B_k=|K'(k)|$, $r_1=1/20$, $r_2=1/25$,
and the fixed $p,c,C$ in \eqref{eq:pressure}, the primitive inner
and outer guards are
\begin{center}
\small
\begin{tabular}{ccc}
\toprule Quantity&$k=1$&$k=2$\\\midrule
$A_k-B_kr_k-2r_k^2$&$36124241/5280000000$&$395330317/154000000000$\\
$q_0=2A_k^2+pk-c$&$45448441/20000000000$&$7255009/125000000000$\\
$B_k-2r_k$&$233921359/264000000$&$2084993683/6160000000$\\
$2Cr_k-p$&$1883/4000$&$1503/4000$\\
$p(k-r_k)+Cr_k^2-c$&$2583/400000$&$259/40000$\\\bottomrule
\end{tabular}
\end{center}
They are all positive. The quadratic coefficients are
\[
 q_2=
 \frac{216248962271606881}{34848000000000000}>6
 \quad(k=1),\qquad
 q_2=\frac{93349805577714704489}{18972800000000000000}>4
 \quad(k=2).
\]
The two exact discriminants are those in
\eqref{eq:discriminants}. The left-interval slopes
$4A_k(B_k-2r_k)-p$ are the positive fractions printed there in
the proof, and $5p/2-c=117/100000>0$ pays every remaining nearest
integer. Thus the rational checks cover the Taylor positivity,
both quadratic minima, the left intervals, outer monotonicity, and
the unbounded range of gaps.

The affine and smoothing arithmetic is
\[
 \begin{gathered}
 Q=1+350(19/4)=3327/2,\\
 B_\star=\frac{350(19/4)B_0+c+6(19/4)}Q
 =\frac{27293111}{332700000},\\
 \frac{350(19/4)270+6}Q=\frac{299254}{1109}.
 \end{gathered}
\]
The pair and span guards can be read without decimal approximation:
\[
 0<a_\star=\frac{725489}{83175000}<\frac9{1000}<\frac1{20}
 <\frac{13707941}{2^{28}},\qquad
 0<b_\star=\frac{584077}{133080000}<\frac{4389}{10^6}.
\]
They also imply $1-a_\star>0$. These are the guards used to pass
from the separated finite certificate to the global count.

\section{Exact data for the twelve-fibre certificate}
\label{app:twelve-fibre}

This appendix specifies all arithmetic used in
\eqref{eq:twelve-certificate-values}: finite rational polynomial
integration and elementary transcendental enclosures.
A numerical boundary-value solve helped find the candidate,
but plays no part in the certificate.

\subsection{The normalized positive candidate}

On cell $j$ put $y=24(s-x_j)-1\in[-1,1]$.
For $0\le j\le5$ define
$p_j(y)=10^{-14}\sum_{k=0}^6n_{jk}y^k$.
The integer data are split into two tables:
\begin{center}
\small
\begin{tabular}{rrrrr}
\toprule
$j$&$n_{j0}$&$n_{j1}$&$n_{j2}$&$n_{j3}$\\
\midrule
0&81452785069661&7292373590683&-495291728588&-8832261778\\
1&94237762244332&5325472480437&-540048842503&-5392125542\\
2&102829051521759&3211889787491&-561973129360&-1752444969\\
3&106611512235993&1220224746606&-555948929227&1548835301\\
4&105872959612158&3026144941748&-520652451189&-4834131202\\
5&110067398521096&987168788008&-542216720202&-1573291099\\
\bottomrule
\end{tabular}
\end{center}
\begin{center}
\small
\begin{tabular}{rrrr}
\toprule
$j$&$n_{j4}$&$n_{j5}$&$n_{j6}$\\
\midrule
0&468777856&3579876&-172336\\
1&498965786&1817587&-181576\\
2&508473653&-36353&-183294\\
3&491383198&-1705511&-175034\\
4&443772398&2367437&-154759\\
5&463005280&777705&-161751\\
\bottomrule
\end{tabular}
\end{center}
Define $p_{11-j}(y)=p_j(-y)$ and divide the full function by
\begin{equation}
 Z_0=\frac{62999999999997821}{63000000000000000}.
 \label{eq:twelve-normalizer}
\end{equation}
Thus $f_0(s)=p_j(y)/Z_0$ on cell $j$.
Using $ds=dy/24$ and
$\int_{-1}^1y^k\dd y=2/(k+1)$ for even $k$, zero for odd $k$,
gives exactly $Z_0$ as the unnormalized mass.
Reflection gives evenness.
On each cell the constant coefficient of $p_j/Z_0$, less the
sum of the absolute values of its other coefficients, exceeds
$7/10$. This proves uniform positivity on the whole cell.

All remaining polynomial operations can be specified by
two linear maps. If $v(y)=\sum_kv_ky^k$, put
\[
 \mathcal I(v)=\frac1{12}
 \sum_{\substack{k\ge0\\k\ {\rm even}}}\frac{v_k}{k+1},
 \qquad
 \mathcal P(v)(y)=\sum_{k\ge0}\frac{v_k}{k+1}
                 \{y^{k+1}-(-1)^{k+1}\}.
\]
Here $\mathcal I$ is the physical-cell integral and
$\mathcal P(v)(-1)=0$.
Write $v_j=p_j/Z_0$. Starting with $\mu_0=0$ and $\nu_0=1/2$,
define
\begin{equation}
 \begin{split}
 P_j(y)&=\mu_j+\mathcal P(v_j)(y)/24,\\
 g_j(y)&=\nu_j+\mathcal P(2P_j-1)(y)/24,\\
 \mu_{j+1}&=P_j(1),\qquad\nu_{j+1}=g_j(1).
 \end{split}
 \label{eq:twelve-potential-recursion}
\end{equation}
The initial potential is $1/2$ by evenness and normalization.
The recursion gives the cumulative integral and exact potential;
$\mu_{12}=1$, $\nu_{12}=1/2$, and
\begin{equation}
 \begin{split}
 A(f_0)&=\sum_j\mathcal I(v_j^2),\\
 C_{k/12}(f_0)&=\sum_{j=0}^{11-k}\mathcal I(v_jv_{j+k}),\\
 J(f_0)&=\sum_j\mathcal I(v_jg_j)
        =2\sum_j\mathcal I(P_j-P_j^2).
 \end{split}
 \label{eq:twelve-exact-integration}
\end{equation}
The last equality supplies a consistency check between two
exact integrations.

\subsection{Seven explicit cosine enclosures}

Every transcendental entry in \eqref{eq:twelve-coefficients}
is one of the seven values below.
The integer $L_r$ in the second column satisfies
\begin{equation}
 \frac{L_r}{10^{18}}<\cos(\pi r)<
 \frac{L_r+1}{10^{18}}.
 \label{eq:twelve-cosine-grid}
\end{equation}
\begin{center}
\small
\begin{tabular}{rr}
\toprule
$r$&$L_r$\\
\midrule
$5189/7500$&$-566923938842617133$\\
$5189/3750$&$-357194495134345020$\\
$10289/20000$&$-45380423447917089$\\
$10289/10000$&$-995881234335375474$\\
$30867/20000$&$135767447683864460$\\
$10289/5000$&$983558865802702074$\\
$30867/10000$&$-963134400298818254$\\
\bottomrule
\end{tabular}
\end{center}
To verify the table, for a positive integer $b$ put
$A_n(b)=\sum_{k=0}^n(-1)^k/((2k+1)b^{2k+1})$.
The alternating remainder and Machin's identity imply
\[
 16A_{65}(5)-4A_{20}(239)
 <\pi<
 16A_{64}(5)-4A_{21}(239).
\]
Reduce each rational multiplier modulo two and by reflection
to $\widetilde r\in[0,1]$.
Multiplication by the displayed interval for $\pi$ gives
$X=[X_-,X_+]$.
Evaluate $\sum_{k=0}^{40}(-1)^kX^{2k}/(2k)!$ using rational
interval arithmetic and enlarge both endpoints by $X_+^{82}/82!$.
Taylor's theorem to degree 81 justifies this enlargement,
since all real derivatives of cosine have absolute value at most one.
The resulting interval lies strictly inside the corresponding
interval in \eqref{eq:twelve-cosine-grid}.
Thus the seven entries follow from the displayed finite sums.

\subsection{Matrix and residual acceptance inequalities}

Insert \eqref{eq:twelve-cosine-grid} into
\eqref{eq:twelve-coefficients}--\eqref{eq:twelve-matrix} with
outward rational interval arithmetic.
The diagonal is exact: $1-qc_0=3/2-17q/12$.
There are three row-diagonal-dominance types.
The last two columns below are strict lower bounds with
denominator $10^{12}$:
\begin{center}
\small
\begin{tabular}{crr}
\toprule
Row indices&$q=q_{25}$&$q=q_{\rm c}$\\
\midrule
$0,1,2,9,10,11$&308234221665&308227513777\\
$3,8$&312599837175&312596684670\\
$4,5,6,7$&355085988964&355079480788\\
\bottomrule
\end{tabular}
\end{center}
Each exceeds $77/250$. The calculation is the exact diagonal
less the sum of the absolute upper bounds of the other entries.

Use \eqref{eq:twelve-exact-integration} in
\eqref{eq:twelve-form} for the energies.
For the residual at $q_{\rm c}$ take
$c=651920180779/10^{12}$ and form on each cell
\[
 r_j(y)=\sum_{k=0}^{11}(M_{q_{\rm c}})_{jk}v_k(y)+g_j(y)-c.
\]
If its coefficient intervals are $[a_{jr},b_{jr}]$, put
$V_j=\sum_r\max(|a_{jr}|,|b_{jr}|)$.
Then $|r_j(y)|\le V_j$ on $[-1,1]$ and
$\|r\|_2^2\le\sum_jV_j^2/12$.
The rational data give
\[
 \begin{gathered}
 651926252579<
 10^{12}\mathcal F_{q_{25}}(f_0)<651926252580,\\
 651920180779<
 10^{12}\mathcal F_{q_{\rm c}}(f_0)<651920180780,\\
 V_j<4/10^{12}\qquad(0\le j\le11).
 \end{gathered}
\]
These are the coarse certificates in
\eqref{eq:twelve-certificate-values}.
Every integral, coefficient and comparison is specified above
as a rational operation on displayed integers or the seven proved
intervals. The residual estimate
\eqref{eq:twelve-residual-certificate} turns these finite
inequalities into the all-$L^2$ conclusion.

\begin{thebibliography}{99}
\raggedright

\bibitem{BGSTB}
Siegfred Alan C. Baluyot, Daniel Alan Goldston, Ade Irma Suriajaya,
and Caroline L. Turnage-Butterbaugh,
\emph{An unconditional Montgomery theorem for pair correlation of zeros
of the Riemann zeta-function},
Acta Arith. \textbf{214} (2024), 357--376.
\url{https://doi.org/10.4064/aa230612-20-3}.
Preprint: \url{https://arxiv.org/abs/2306.04799v1}.

\bibitem{AF}
Levent Alp\"oge and Ralph Furman,
\emph{More than two thirds of the zeta zeros are simple and on the critical line},
preprint, arXiv:2608.13637v2, 19 August 2026.
\url{https://arxiv.org/abs/2608.13637v2}.

\bibitem{Lamzouri}
Youness Lamzouri,
\emph{A new proof that more than 2/3 of the zeros of the Riemann zeta
function are simple and on the critical line},
preprint, arXiv:2609.02882v2, 8 September 2026.
\url{https://arxiv.org/abs/2609.02882v2}.

\bibitem{ChandruP6}
R. Arun Chandru,
\emph{Compensated Gram Corrections for Zero Counts of the Riemann Zeta Function},
Zenodo preprint, version\_1, 11 September 2026
(manuscript dated 10 September 2026).
\url{https://doi.org/10.5281/zenodo.22701203}.

\bibitem{Ainta}
ainta,
\emph{More than 67.3\% of the zeros of the Riemann zeta function
are simple and lie on the critical line},
public research repository, pinned commit of 11 August 2026.
\url{https://github.com/ainta/zeta-simple-zeros/tree/040c5e899e658aed7b56a2a87f501798fe10761d}.

\bibitem{AMTOPA}
AMTOPA,
\emph{Joint window and exact-pressure optimization for simple zeros
of the Riemann zeta function},
public research repository, pinned commit of 13 August 2026.
\url{https://github.com/AMTOPA/zeta-exact-pressure/tree/7253fdcab9366af45b8c8caf44e408c0af44a1a7}.

\bibitem{GebendorferWindows}
Jonas Jakob Gebendorfer,
\emph{Compatible local windows and a certified lower bound
for simple critical-line zeros},
research draft dated 7 September 2026; Zenodo deposit, 8 September 2026.
\url{https://doi.org/10.5281/zenodo.22661276}.

\bibitem{GebendorferExpanded}
Jonas Jakob Gebendorfer,
\emph{An expanded window and joint position pressures
for simple critical-line zeros},
research draft v0.3, 7 September 2026.
\url{https://doi.org/10.5281/zenodo.22643957}.

\bibitem{Tawanerguo}
tawanerguo-cn,
\emph{A certified unconditional 67.3192911473\% lower bound
for simple zeros of the Riemann zeta function},
public research repository, pinned commit of 11 August 2026;
archive release dated 12 August 2026.
\url{https://github.com/tawanerguo-cn/zeta-simple-zeros/tree/45149f6d403059a71be73c5e3f884cee7cd62b20}.
Archive: \url{https://doi.org/10.5281/zenodo.21892048}.

\bibitem{Knausgard2026}
Kristian Muri Knausg{\aa}rd,
\emph{More than 83.9\% of the zeros of the Riemann zeta function
are distinct and more than 67.35\% are simple and on the critical line},
preprint, arXiv:2610.08965v1, 6 October 2026.
\url{https://arxiv.org/abs/2610.08965v1}.

\bibitem{Freund}
G. Freund,
\emph{Lossless local-pressure transfer by separating close pairs},
draft pull request 888, 14 September 2026.
\url{https://github.com/GettysburgResearch/riemann/pull/888}.
Pinned proof:
\url{https://github.com/GettysburgResearch/riemann/blob/1ce38f1f5b6f39b5a182261c33936c23028e40ea/standalone/2026-09-14-lossless-pressure/PROOF.md}.

\bibitem{DSouza}
Oliver D'Souza,
\emph{On the proportions of simple and of distinct zeros
of the Riemann zeta function},
research preprint with Lean formalization, release v1.0, 3 October 2026.
\url{https://github.com/Oliverds321/simple-zeros/tree/c8fee8e21ad93b3504b215d6a6098b96cb7a0cea}.

\bibitem{GebendorferCap}
Jonas Jakob Gebendorfer,
\emph{Compatible triples and a periodic-mixture obstruction
for a critical-line zero certificate},
consolidated research draft, 7 September 2026.
\url{https://doi.org/10.5281/zenodo.22647797}.

\bibitem{OjhaTrmdy}
Vivaswat Ojha,
\emph{A certified 67.32001\% lower-bound candidate for simple zeros
of the Riemann zeta function on the critical line},
contributor draft dated 11 August 2026, in the public research repository
of trmdy, pinned at
\url{https://github.com/trmdy/zeta-simple-zeros-673137/tree/1610b97b7895ff34982260f8dcaf04a0f7b82cf7}.
The related repository report
\emph{Campaign 2 (2026-08-11, afternoon): three attempted improvements
and what they proved} is at
\url{https://github.com/trmdy/zeta-simple-zeros-673137/blob/1610b97b7895ff34982260f8dcaf04a0f7b82cf7/docs/campaign-2.md}.

\bibitem{Rakitin}
Dmitry Rakitin,
\emph{Zeta and Xi Zero Results},
public research repository, release v0.1.1, 15 September 2026.
\url{https://github.com/RakitinD/zeta-xi-zero-results/tree/8dd1ba8db18c3d27effcf3257577587513b6cd98}.

\bibitem{ShiTrace}
Yuhang Shi,
\emph{A two-certificate trace-energy deduction for simple zeros
of the Riemann zeta function},
public research draft v0.1.0, 14 August 2026.
\url{https://doi.org/10.5281/zenodo.21926962}.

\bibitem{Schafer}
Dirk Sch\"afer,
\emph{A Certified Finite Refinement for Simple and Distinct Zeros
of the Riemann Zeta Function},
preprint, 18 September 2026.
\url{https://doi.org/10.5281/zenodo.22837912}.

\bibitem{Singh}
Ketan Singh,
\emph{An Unconditional 67.345129\% Lower Bound for Simple Zeros
of the Riemann Zeta Function on the Critical Line},
preprint v1, 30 August 2026.
\url{https://doi.org/10.5281/zenodo.22119371}.

\bibitem{KnausgardDistinct}
Kristian Muri Knausg{\aa}rd,
\emph{More than 83.69\% of the zeros of the Riemann zeta function
are distinct},
preprint, arXiv:2609.33043v2, 7 October 2026.
\url{https://arxiv.org/abs/2609.33043v2}.



\bibitem{BGSTBI}
Siegfred Alan C. Baluyot, Daniel Alan Goldston, Ade Irma Suriajaya,
and Caroline L. Turnage-Butterbaugh,
\emph{Pair Correlation of Zeros of the Riemann Zeta Function I:
Proportions of Simple Zeros and Critical Zeros},
preprint, arXiv:2501.14545v3, 1 September 2026.
\url{https://arxiv.org/abs/2501.14545v3}.

\bibitem{Montgomery}
Hugh L. Montgomery,
\emph{The pair correlation of zeros of the zeta function},
Proc. Sympos. Pure Math. \textbf{24}, American Mathematical Society,
1973, 181--193.

\bibitem{MontgomeryTaylor}
Hugh L. Montgomery,
\emph{Distribution of the zeros of the Riemann zeta function},
Proceedings of the International Congress of Mathematicians
(Vancouver, 1974), vol.~1, Canadian Mathematical Congress,
1975, 379--381.

\bibitem{CGdL}
Andr\'es Chirre, Felipe Gon\c{c}alves and David de Laat,
\emph{Pair correlation estimates for the zeros of the zeta function
via semidefinite programming},
Adv. Math. \textbf{361} (2020), 106926.
\url{https://arxiv.org/abs/1810.08843v2}.

\bibitem{Santibanez}
Felipe Santiba\~nez-Leal,
\emph{Stronger local overlap certificates and distinct zeros of the
Riemann zeta function}, preprint, version 0.01, 3 October 2026.
\url{https://doi.org/10.5281/zenodo.23128663}.

\bibitem{Lavery2026}
Samuel Lavery,
\emph{Eight-point weighted refinement with edge windows},
Lean source and certificate data, 2026.
\url{https://github.com/josusanmartin/riemann/blob/296f7bb9c4e01b0c97554cf76af91ed99a410e8d/submissions/attempt-016/proof/Solution.lean}.

\begingroup\interlinepenalty=10000
\bibitem{Waddle}
Zach Waddle,
\emph{Unconditional low-multiplicity profiles for zeros of the Riemann
zeta function: Sharp consequences of a multiplicity-aware rank--trace
inequality}, public repository draft, version 0.2, 12 August 2026.
\url{https://github.com/zach7036/riemann-hypothesis-research/blob/main/publication/low-multiplicity-zeta/low_multiplicity_zeta.tex}.
\par\endgroup

\bibitem{STT}
Aleksander Simoni\v{c}, Timothy S. Trudgian and Caroline L. Turnage-Butterbaugh,
\emph{Some explicit and unconditional results on gaps between zeroes
of the Riemann zeta-function},
Trans. Amer. Math. Soc. \textbf{375} (2022), 3239--3265.
\url{https://doi.org/10.1090/tran/8571}.

\bibitem{SM}
Accompanying Supplemental Material is in the repository - 
\href{https://github.com/arun-chandru/riemann-zeros}{github.com/arun-chandru/riemann-zeros}

\end{thebibliography}
\end{document}



